\section*{Item 3}\label{it:3}

\textbf{Enunciado}: para um conversor Cuk, considere os seguintes valores: $E = 48$ V; $V_o = 36$ V; $R_o = 9~\Omega$; $f = 64$ kHz; $L_1 = 10$ mH; $L2 = 1$ mH; $C_o = 100~\mu$F; rendimento de 100\%.

\medskip
\textbf{a)} Determine se o conversor está operando no CCM ou no DCM.

\medskip
\textbf{Resolução}: pode-se verificar o modo de operação do conversor por meio do limiar inferior da corrente~\cite{aula3.3}
\begin{equation*}
	\begin{aligned}
		I_{o,min} &= \frac{ET_s}{8L_2} \\
		&= \frac{48(1/64\times 10^3)}{8(1 \times 10^{-3})} \approx 0{,}0938~\mathrm{A}.
	\end{aligned}
\end{equation*}

Comparando-se $I_{o,min}$ com $I_{o} = V_o/R_o = 36/9 = 4$ A, conclui-se que o conversor está operando em CCM. Essa condição pode ser verificada por meio de simulação da corrente do diodo, como mostrado na Figura~\ref{fig:ex3_1}.
\begin{figure}[htb!]
	\centering
	%% Creator: Matplotlib, PGF backend
%%
%% To include the figure in your LaTeX document, write
%%   \input{<filename>.pgf}
%%
%% Make sure the required packages are loaded in your preamble
%%   \usepackage{pgf}
%%
%% Also ensure that all the required font packages are loaded; for instance,
%% the lmodern package is sometimes necessary when using math font.
%%   \usepackage{lmodern}
%%
%% Figures using additional raster images can only be included by \input if
%% they are in the same directory as the main LaTeX file. For loading figures
%% from other directories you can use the `import` package
%%   \usepackage{import}
%%
%% and then include the figures with
%%   \import{<path to file>}{<filename>.pgf}
%%
%% Matplotlib used the following preamble
%%   \def\mathdefault#1{#1}
%%   \everymath=\expandafter{\the\everymath\displaystyle}
%%   \IfFileExists{scrextend.sty}{
%%     \usepackage[fontsize=8.000000pt]{scrextend}
%%   }{
%%     \renewcommand{\normalsize}{\fontsize{8.000000}{9.600000}\selectfont}
%%     \normalsize
%%   }
%%   
%%           \usepackage{amsmath}
%%           \usepackage{amssymb}
%%           \usepackage{mathtools}
%%           \usepackage{dsfont}
%%           \providecommand{\mathdefault}[1]{#1}
%%       
%%   \makeatletter\@ifpackageloaded{underscore}{}{\usepackage[strings]{underscore}}\makeatother
%%
\begingroup%
\makeatletter%
\begin{pgfpicture}%
\pgfpathrectangle{\pgfpointorigin}{\pgfqpoint{3.242918in}{1.514654in}}%
\pgfusepath{use as bounding box, clip}%
\begin{pgfscope}%
\pgfsetbuttcap%
\pgfsetmiterjoin%
\definecolor{currentfill}{rgb}{1.000000,1.000000,1.000000}%
\pgfsetfillcolor{currentfill}%
\pgfsetlinewidth{0.000000pt}%
\definecolor{currentstroke}{rgb}{1.000000,1.000000,1.000000}%
\pgfsetstrokecolor{currentstroke}%
\pgfsetdash{}{0pt}%
\pgfpathmoveto{\pgfqpoint{0.000000in}{0.000000in}}%
\pgfpathlineto{\pgfqpoint{3.242918in}{0.000000in}}%
\pgfpathlineto{\pgfqpoint{3.242918in}{1.514654in}}%
\pgfpathlineto{\pgfqpoint{0.000000in}{1.514654in}}%
\pgfpathlineto{\pgfqpoint{0.000000in}{0.000000in}}%
\pgfpathclose%
\pgfusepath{fill}%
\end{pgfscope}%
\begin{pgfscope}%
\pgfsetbuttcap%
\pgfsetmiterjoin%
\definecolor{currentfill}{rgb}{1.000000,1.000000,1.000000}%
\pgfsetfillcolor{currentfill}%
\pgfsetlinewidth{0.000000pt}%
\definecolor{currentstroke}{rgb}{0.000000,0.000000,0.000000}%
\pgfsetstrokecolor{currentstroke}%
\pgfsetstrokeopacity{0.000000}%
\pgfsetdash{}{0pt}%
\pgfpathmoveto{\pgfqpoint{0.422918in}{0.462654in}}%
\pgfpathlineto{\pgfqpoint{3.142917in}{0.462654in}}%
\pgfpathlineto{\pgfqpoint{3.142917in}{1.414654in}}%
\pgfpathlineto{\pgfqpoint{0.422918in}{1.414654in}}%
\pgfpathlineto{\pgfqpoint{0.422918in}{0.462654in}}%
\pgfpathclose%
\pgfusepath{fill}%
\end{pgfscope}%
\begin{pgfscope}%
\pgfpathrectangle{\pgfqpoint{0.422918in}{0.462654in}}{\pgfqpoint{2.720000in}{0.952000in}}%
\pgfusepath{clip}%
\pgfsetbuttcap%
\pgfsetroundjoin%
\pgfsetlinewidth{0.501875pt}%
\definecolor{currentstroke}{rgb}{0.690196,0.690196,0.690196}%
\pgfsetstrokecolor{currentstroke}%
\pgfsetstrokeopacity{0.250000}%
\pgfsetdash{{1.850000pt}{0.800000pt}}{0.000000pt}%
\pgfpathmoveto{\pgfqpoint{0.777342in}{0.462654in}}%
\pgfpathlineto{\pgfqpoint{0.777342in}{1.414654in}}%
\pgfusepath{stroke}%
\end{pgfscope}%
\begin{pgfscope}%
\pgfsetbuttcap%
\pgfsetroundjoin%
\definecolor{currentfill}{rgb}{0.000000,0.000000,0.000000}%
\pgfsetfillcolor{currentfill}%
\pgfsetlinewidth{0.803000pt}%
\definecolor{currentstroke}{rgb}{0.000000,0.000000,0.000000}%
\pgfsetstrokecolor{currentstroke}%
\pgfsetdash{}{0pt}%
\pgfsys@defobject{currentmarker}{\pgfqpoint{0.000000in}{-0.048611in}}{\pgfqpoint{0.000000in}{0.000000in}}{%
\pgfpathmoveto{\pgfqpoint{0.000000in}{0.000000in}}%
\pgfpathlineto{\pgfqpoint{0.000000in}{-0.048611in}}%
\pgfusepath{stroke,fill}%
}%
\begin{pgfscope}%
\pgfsys@transformshift{0.777342in}{0.462654in}%
\pgfsys@useobject{currentmarker}{}%
\end{pgfscope}%
\end{pgfscope}%
\begin{pgfscope}%
\definecolor{textcolor}{rgb}{0.000000,0.000000,0.000000}%
\pgfsetstrokecolor{textcolor}%
\pgfsetfillcolor{textcolor}%
\pgftext[x=0.777342in,y=0.365432in,,top]{\color{textcolor}{\rmfamily\fontsize{8.000000}{9.600000}\selectfont\catcode`\^=\active\def^{\ifmmode\sp\else\^{}\fi}\catcode`\%=\active\def%{\%}$\mathdefault{7.804}$}}%
\end{pgfscope}%
\begin{pgfscope}%
\pgfpathrectangle{\pgfqpoint{0.422918in}{0.462654in}}{\pgfqpoint{2.720000in}{0.952000in}}%
\pgfusepath{clip}%
\pgfsetbuttcap%
\pgfsetroundjoin%
\pgfsetlinewidth{0.501875pt}%
\definecolor{currentstroke}{rgb}{0.690196,0.690196,0.690196}%
\pgfsetstrokecolor{currentstroke}%
\pgfsetstrokeopacity{0.250000}%
\pgfsetdash{{1.850000pt}{0.800000pt}}{0.000000pt}%
\pgfpathmoveto{\pgfqpoint{1.304857in}{0.462654in}}%
\pgfpathlineto{\pgfqpoint{1.304857in}{1.414654in}}%
\pgfusepath{stroke}%
\end{pgfscope}%
\begin{pgfscope}%
\pgfsetbuttcap%
\pgfsetroundjoin%
\definecolor{currentfill}{rgb}{0.000000,0.000000,0.000000}%
\pgfsetfillcolor{currentfill}%
\pgfsetlinewidth{0.803000pt}%
\definecolor{currentstroke}{rgb}{0.000000,0.000000,0.000000}%
\pgfsetstrokecolor{currentstroke}%
\pgfsetdash{}{0pt}%
\pgfsys@defobject{currentmarker}{\pgfqpoint{0.000000in}{-0.048611in}}{\pgfqpoint{0.000000in}{0.000000in}}{%
\pgfpathmoveto{\pgfqpoint{0.000000in}{0.000000in}}%
\pgfpathlineto{\pgfqpoint{0.000000in}{-0.048611in}}%
\pgfusepath{stroke,fill}%
}%
\begin{pgfscope}%
\pgfsys@transformshift{1.304857in}{0.462654in}%
\pgfsys@useobject{currentmarker}{}%
\end{pgfscope}%
\end{pgfscope}%
\begin{pgfscope}%
\definecolor{textcolor}{rgb}{0.000000,0.000000,0.000000}%
\pgfsetstrokecolor{textcolor}%
\pgfsetfillcolor{textcolor}%
\pgftext[x=1.304857in,y=0.365432in,,top]{\color{textcolor}{\rmfamily\fontsize{8.000000}{9.600000}\selectfont\catcode`\^=\active\def^{\ifmmode\sp\else\^{}\fi}\catcode`\%=\active\def%{\%}$\mathdefault{7.806}$}}%
\end{pgfscope}%
\begin{pgfscope}%
\pgfpathrectangle{\pgfqpoint{0.422918in}{0.462654in}}{\pgfqpoint{2.720000in}{0.952000in}}%
\pgfusepath{clip}%
\pgfsetbuttcap%
\pgfsetroundjoin%
\pgfsetlinewidth{0.501875pt}%
\definecolor{currentstroke}{rgb}{0.690196,0.690196,0.690196}%
\pgfsetstrokecolor{currentstroke}%
\pgfsetstrokeopacity{0.250000}%
\pgfsetdash{{1.850000pt}{0.800000pt}}{0.000000pt}%
\pgfpathmoveto{\pgfqpoint{1.832372in}{0.462654in}}%
\pgfpathlineto{\pgfqpoint{1.832372in}{1.414654in}}%
\pgfusepath{stroke}%
\end{pgfscope}%
\begin{pgfscope}%
\pgfsetbuttcap%
\pgfsetroundjoin%
\definecolor{currentfill}{rgb}{0.000000,0.000000,0.000000}%
\pgfsetfillcolor{currentfill}%
\pgfsetlinewidth{0.803000pt}%
\definecolor{currentstroke}{rgb}{0.000000,0.000000,0.000000}%
\pgfsetstrokecolor{currentstroke}%
\pgfsetdash{}{0pt}%
\pgfsys@defobject{currentmarker}{\pgfqpoint{0.000000in}{-0.048611in}}{\pgfqpoint{0.000000in}{0.000000in}}{%
\pgfpathmoveto{\pgfqpoint{0.000000in}{0.000000in}}%
\pgfpathlineto{\pgfqpoint{0.000000in}{-0.048611in}}%
\pgfusepath{stroke,fill}%
}%
\begin{pgfscope}%
\pgfsys@transformshift{1.832372in}{0.462654in}%
\pgfsys@useobject{currentmarker}{}%
\end{pgfscope}%
\end{pgfscope}%
\begin{pgfscope}%
\definecolor{textcolor}{rgb}{0.000000,0.000000,0.000000}%
\pgfsetstrokecolor{textcolor}%
\pgfsetfillcolor{textcolor}%
\pgftext[x=1.832372in,y=0.365432in,,top]{\color{textcolor}{\rmfamily\fontsize{8.000000}{9.600000}\selectfont\catcode`\^=\active\def^{\ifmmode\sp\else\^{}\fi}\catcode`\%=\active\def%{\%}$\mathdefault{7.808}$}}%
\end{pgfscope}%
\begin{pgfscope}%
\pgfpathrectangle{\pgfqpoint{0.422918in}{0.462654in}}{\pgfqpoint{2.720000in}{0.952000in}}%
\pgfusepath{clip}%
\pgfsetbuttcap%
\pgfsetroundjoin%
\pgfsetlinewidth{0.501875pt}%
\definecolor{currentstroke}{rgb}{0.690196,0.690196,0.690196}%
\pgfsetstrokecolor{currentstroke}%
\pgfsetstrokeopacity{0.250000}%
\pgfsetdash{{1.850000pt}{0.800000pt}}{0.000000pt}%
\pgfpathmoveto{\pgfqpoint{2.359887in}{0.462654in}}%
\pgfpathlineto{\pgfqpoint{2.359887in}{1.414654in}}%
\pgfusepath{stroke}%
\end{pgfscope}%
\begin{pgfscope}%
\pgfsetbuttcap%
\pgfsetroundjoin%
\definecolor{currentfill}{rgb}{0.000000,0.000000,0.000000}%
\pgfsetfillcolor{currentfill}%
\pgfsetlinewidth{0.803000pt}%
\definecolor{currentstroke}{rgb}{0.000000,0.000000,0.000000}%
\pgfsetstrokecolor{currentstroke}%
\pgfsetdash{}{0pt}%
\pgfsys@defobject{currentmarker}{\pgfqpoint{0.000000in}{-0.048611in}}{\pgfqpoint{0.000000in}{0.000000in}}{%
\pgfpathmoveto{\pgfqpoint{0.000000in}{0.000000in}}%
\pgfpathlineto{\pgfqpoint{0.000000in}{-0.048611in}}%
\pgfusepath{stroke,fill}%
}%
\begin{pgfscope}%
\pgfsys@transformshift{2.359887in}{0.462654in}%
\pgfsys@useobject{currentmarker}{}%
\end{pgfscope}%
\end{pgfscope}%
\begin{pgfscope}%
\definecolor{textcolor}{rgb}{0.000000,0.000000,0.000000}%
\pgfsetstrokecolor{textcolor}%
\pgfsetfillcolor{textcolor}%
\pgftext[x=2.359887in,y=0.365432in,,top]{\color{textcolor}{\rmfamily\fontsize{8.000000}{9.600000}\selectfont\catcode`\^=\active\def^{\ifmmode\sp\else\^{}\fi}\catcode`\%=\active\def%{\%}$\mathdefault{7.810}$}}%
\end{pgfscope}%
\begin{pgfscope}%
\pgfpathrectangle{\pgfqpoint{0.422918in}{0.462654in}}{\pgfqpoint{2.720000in}{0.952000in}}%
\pgfusepath{clip}%
\pgfsetbuttcap%
\pgfsetroundjoin%
\pgfsetlinewidth{0.501875pt}%
\definecolor{currentstroke}{rgb}{0.690196,0.690196,0.690196}%
\pgfsetstrokecolor{currentstroke}%
\pgfsetstrokeopacity{0.250000}%
\pgfsetdash{{1.850000pt}{0.800000pt}}{0.000000pt}%
\pgfpathmoveto{\pgfqpoint{2.887402in}{0.462654in}}%
\pgfpathlineto{\pgfqpoint{2.887402in}{1.414654in}}%
\pgfusepath{stroke}%
\end{pgfscope}%
\begin{pgfscope}%
\pgfsetbuttcap%
\pgfsetroundjoin%
\definecolor{currentfill}{rgb}{0.000000,0.000000,0.000000}%
\pgfsetfillcolor{currentfill}%
\pgfsetlinewidth{0.803000pt}%
\definecolor{currentstroke}{rgb}{0.000000,0.000000,0.000000}%
\pgfsetstrokecolor{currentstroke}%
\pgfsetdash{}{0pt}%
\pgfsys@defobject{currentmarker}{\pgfqpoint{0.000000in}{-0.048611in}}{\pgfqpoint{0.000000in}{0.000000in}}{%
\pgfpathmoveto{\pgfqpoint{0.000000in}{0.000000in}}%
\pgfpathlineto{\pgfqpoint{0.000000in}{-0.048611in}}%
\pgfusepath{stroke,fill}%
}%
\begin{pgfscope}%
\pgfsys@transformshift{2.887402in}{0.462654in}%
\pgfsys@useobject{currentmarker}{}%
\end{pgfscope}%
\end{pgfscope}%
\begin{pgfscope}%
\definecolor{textcolor}{rgb}{0.000000,0.000000,0.000000}%
\pgfsetstrokecolor{textcolor}%
\pgfsetfillcolor{textcolor}%
\pgftext[x=2.887402in,y=0.365432in,,top]{\color{textcolor}{\rmfamily\fontsize{8.000000}{9.600000}\selectfont\catcode`\^=\active\def^{\ifmmode\sp\else\^{}\fi}\catcode`\%=\active\def%{\%}$\mathdefault{7.812}$}}%
\end{pgfscope}%
\begin{pgfscope}%
\definecolor{textcolor}{rgb}{0.000000,0.000000,0.000000}%
\pgfsetstrokecolor{textcolor}%
\pgfsetfillcolor{textcolor}%
\pgftext[x=1.782917in,y=0.211111in,,top]{\color{textcolor}{\rmfamily\fontsize{8.000000}{9.600000}\selectfont\catcode`\^=\active\def^{\ifmmode\sp\else\^{}\fi}\catcode`\%=\active\def%{\%}Time (ms)}}%
\end{pgfscope}%
\begin{pgfscope}%
\pgfpathrectangle{\pgfqpoint{0.422918in}{0.462654in}}{\pgfqpoint{2.720000in}{0.952000in}}%
\pgfusepath{clip}%
\pgfsetbuttcap%
\pgfsetroundjoin%
\pgfsetlinewidth{0.501875pt}%
\definecolor{currentstroke}{rgb}{0.690196,0.690196,0.690196}%
\pgfsetstrokecolor{currentstroke}%
\pgfsetstrokeopacity{0.250000}%
\pgfsetdash{{1.850000pt}{0.800000pt}}{0.000000pt}%
\pgfpathmoveto{\pgfqpoint{0.422918in}{0.506048in}}%
\pgfpathlineto{\pgfqpoint{3.142917in}{0.506048in}}%
\pgfusepath{stroke}%
\end{pgfscope}%
\begin{pgfscope}%
\pgfsetbuttcap%
\pgfsetroundjoin%
\definecolor{currentfill}{rgb}{0.000000,0.000000,0.000000}%
\pgfsetfillcolor{currentfill}%
\pgfsetlinewidth{0.803000pt}%
\definecolor{currentstroke}{rgb}{0.000000,0.000000,0.000000}%
\pgfsetstrokecolor{currentstroke}%
\pgfsetdash{}{0pt}%
\pgfsys@defobject{currentmarker}{\pgfqpoint{-0.048611in}{0.000000in}}{\pgfqpoint{-0.000000in}{0.000000in}}{%
\pgfpathmoveto{\pgfqpoint{-0.000000in}{0.000000in}}%
\pgfpathlineto{\pgfqpoint{-0.048611in}{0.000000in}}%
\pgfusepath{stroke,fill}%
}%
\begin{pgfscope}%
\pgfsys@transformshift{0.422918in}{0.506048in}%
\pgfsys@useobject{currentmarker}{}%
\end{pgfscope}%
\end{pgfscope}%
\begin{pgfscope}%
\definecolor{textcolor}{rgb}{0.000000,0.000000,0.000000}%
\pgfsetstrokecolor{textcolor}%
\pgfsetfillcolor{textcolor}%
\pgftext[x=0.266667in, y=0.467468in, left, base]{\color{textcolor}{\rmfamily\fontsize{8.000000}{9.600000}\selectfont\catcode`\^=\active\def^{\ifmmode\sp\else\^{}\fi}\catcode`\%=\active\def%{\%}$\mathdefault{0}$}}%
\end{pgfscope}%
\begin{pgfscope}%
\pgfpathrectangle{\pgfqpoint{0.422918in}{0.462654in}}{\pgfqpoint{2.720000in}{0.952000in}}%
\pgfusepath{clip}%
\pgfsetbuttcap%
\pgfsetroundjoin%
\pgfsetlinewidth{0.501875pt}%
\definecolor{currentstroke}{rgb}{0.690196,0.690196,0.690196}%
\pgfsetstrokecolor{currentstroke}%
\pgfsetstrokeopacity{0.250000}%
\pgfsetdash{{1.850000pt}{0.800000pt}}{0.000000pt}%
\pgfpathmoveto{\pgfqpoint{0.422918in}{0.748082in}}%
\pgfpathlineto{\pgfqpoint{3.142917in}{0.748082in}}%
\pgfusepath{stroke}%
\end{pgfscope}%
\begin{pgfscope}%
\pgfsetbuttcap%
\pgfsetroundjoin%
\definecolor{currentfill}{rgb}{0.000000,0.000000,0.000000}%
\pgfsetfillcolor{currentfill}%
\pgfsetlinewidth{0.803000pt}%
\definecolor{currentstroke}{rgb}{0.000000,0.000000,0.000000}%
\pgfsetstrokecolor{currentstroke}%
\pgfsetdash{}{0pt}%
\pgfsys@defobject{currentmarker}{\pgfqpoint{-0.048611in}{0.000000in}}{\pgfqpoint{-0.000000in}{0.000000in}}{%
\pgfpathmoveto{\pgfqpoint{-0.000000in}{0.000000in}}%
\pgfpathlineto{\pgfqpoint{-0.048611in}{0.000000in}}%
\pgfusepath{stroke,fill}%
}%
\begin{pgfscope}%
\pgfsys@transformshift{0.422918in}{0.748082in}%
\pgfsys@useobject{currentmarker}{}%
\end{pgfscope}%
\end{pgfscope}%
\begin{pgfscope}%
\definecolor{textcolor}{rgb}{0.000000,0.000000,0.000000}%
\pgfsetstrokecolor{textcolor}%
\pgfsetfillcolor{textcolor}%
\pgftext[x=0.266667in, y=0.709502in, left, base]{\color{textcolor}{\rmfamily\fontsize{8.000000}{9.600000}\selectfont\catcode`\^=\active\def^{\ifmmode\sp\else\^{}\fi}\catcode`\%=\active\def%{\%}$\mathdefault{2}$}}%
\end{pgfscope}%
\begin{pgfscope}%
\pgfpathrectangle{\pgfqpoint{0.422918in}{0.462654in}}{\pgfqpoint{2.720000in}{0.952000in}}%
\pgfusepath{clip}%
\pgfsetbuttcap%
\pgfsetroundjoin%
\pgfsetlinewidth{0.501875pt}%
\definecolor{currentstroke}{rgb}{0.690196,0.690196,0.690196}%
\pgfsetstrokecolor{currentstroke}%
\pgfsetstrokeopacity{0.250000}%
\pgfsetdash{{1.850000pt}{0.800000pt}}{0.000000pt}%
\pgfpathmoveto{\pgfqpoint{0.422918in}{0.990115in}}%
\pgfpathlineto{\pgfqpoint{3.142917in}{0.990115in}}%
\pgfusepath{stroke}%
\end{pgfscope}%
\begin{pgfscope}%
\pgfsetbuttcap%
\pgfsetroundjoin%
\definecolor{currentfill}{rgb}{0.000000,0.000000,0.000000}%
\pgfsetfillcolor{currentfill}%
\pgfsetlinewidth{0.803000pt}%
\definecolor{currentstroke}{rgb}{0.000000,0.000000,0.000000}%
\pgfsetstrokecolor{currentstroke}%
\pgfsetdash{}{0pt}%
\pgfsys@defobject{currentmarker}{\pgfqpoint{-0.048611in}{0.000000in}}{\pgfqpoint{-0.000000in}{0.000000in}}{%
\pgfpathmoveto{\pgfqpoint{-0.000000in}{0.000000in}}%
\pgfpathlineto{\pgfqpoint{-0.048611in}{0.000000in}}%
\pgfusepath{stroke,fill}%
}%
\begin{pgfscope}%
\pgfsys@transformshift{0.422918in}{0.990115in}%
\pgfsys@useobject{currentmarker}{}%
\end{pgfscope}%
\end{pgfscope}%
\begin{pgfscope}%
\definecolor{textcolor}{rgb}{0.000000,0.000000,0.000000}%
\pgfsetstrokecolor{textcolor}%
\pgfsetfillcolor{textcolor}%
\pgftext[x=0.266667in, y=0.951535in, left, base]{\color{textcolor}{\rmfamily\fontsize{8.000000}{9.600000}\selectfont\catcode`\^=\active\def^{\ifmmode\sp\else\^{}\fi}\catcode`\%=\active\def%{\%}$\mathdefault{4}$}}%
\end{pgfscope}%
\begin{pgfscope}%
\pgfpathrectangle{\pgfqpoint{0.422918in}{0.462654in}}{\pgfqpoint{2.720000in}{0.952000in}}%
\pgfusepath{clip}%
\pgfsetbuttcap%
\pgfsetroundjoin%
\pgfsetlinewidth{0.501875pt}%
\definecolor{currentstroke}{rgb}{0.690196,0.690196,0.690196}%
\pgfsetstrokecolor{currentstroke}%
\pgfsetstrokeopacity{0.250000}%
\pgfsetdash{{1.850000pt}{0.800000pt}}{0.000000pt}%
\pgfpathmoveto{\pgfqpoint{0.422918in}{1.232149in}}%
\pgfpathlineto{\pgfqpoint{3.142917in}{1.232149in}}%
\pgfusepath{stroke}%
\end{pgfscope}%
\begin{pgfscope}%
\pgfsetbuttcap%
\pgfsetroundjoin%
\definecolor{currentfill}{rgb}{0.000000,0.000000,0.000000}%
\pgfsetfillcolor{currentfill}%
\pgfsetlinewidth{0.803000pt}%
\definecolor{currentstroke}{rgb}{0.000000,0.000000,0.000000}%
\pgfsetstrokecolor{currentstroke}%
\pgfsetdash{}{0pt}%
\pgfsys@defobject{currentmarker}{\pgfqpoint{-0.048611in}{0.000000in}}{\pgfqpoint{-0.000000in}{0.000000in}}{%
\pgfpathmoveto{\pgfqpoint{-0.000000in}{0.000000in}}%
\pgfpathlineto{\pgfqpoint{-0.048611in}{0.000000in}}%
\pgfusepath{stroke,fill}%
}%
\begin{pgfscope}%
\pgfsys@transformshift{0.422918in}{1.232149in}%
\pgfsys@useobject{currentmarker}{}%
\end{pgfscope}%
\end{pgfscope}%
\begin{pgfscope}%
\definecolor{textcolor}{rgb}{0.000000,0.000000,0.000000}%
\pgfsetstrokecolor{textcolor}%
\pgfsetfillcolor{textcolor}%
\pgftext[x=0.266667in, y=1.193569in, left, base]{\color{textcolor}{\rmfamily\fontsize{8.000000}{9.600000}\selectfont\catcode`\^=\active\def^{\ifmmode\sp\else\^{}\fi}\catcode`\%=\active\def%{\%}$\mathdefault{6}$}}%
\end{pgfscope}%
\begin{pgfscope}%
\definecolor{textcolor}{rgb}{0.000000,0.000000,0.000000}%
\pgfsetstrokecolor{textcolor}%
\pgfsetfillcolor{textcolor}%
\pgftext[x=0.211111in,y=0.938654in,,bottom,rotate=90.000000]{\color{textcolor}{\rmfamily\fontsize{8.000000}{9.600000}\selectfont\catcode`\^=\active\def^{\ifmmode\sp\else\^{}\fi}\catcode`\%=\active\def%{\%}$i_D$ (A)}}%
\end{pgfscope}%
\begin{pgfscope}%
\pgfpathrectangle{\pgfqpoint{0.422918in}{0.462654in}}{\pgfqpoint{2.720000in}{0.952000in}}%
\pgfusepath{clip}%
\pgfsetrectcap%
\pgfsetroundjoin%
\pgfsetlinewidth{1.505625pt}%
\definecolor{currentstroke}{rgb}{0.835294,0.368627,0.000000}%
\pgfsetstrokecolor{currentstroke}%
\pgfsetdash{}{0pt}%
\pgfpathmoveto{\pgfqpoint{0.546554in}{1.326907in}}%
\pgfpathlineto{\pgfqpoint{0.546564in}{1.326904in}}%
\pgfpathlineto{\pgfqpoint{0.547102in}{0.505927in}}%
\pgfpathlineto{\pgfqpoint{0.723224in}{0.505927in}}%
\pgfpathlineto{\pgfqpoint{0.723229in}{1.371382in}}%
\pgfpathlineto{\pgfqpoint{0.723915in}{1.369579in}}%
\pgfpathlineto{\pgfqpoint{0.958686in}{1.326904in}}%
\pgfpathlineto{\pgfqpoint{0.959224in}{0.505927in}}%
\pgfpathlineto{\pgfqpoint{1.135345in}{0.505927in}}%
\pgfpathlineto{\pgfqpoint{1.135350in}{1.371382in}}%
\pgfpathlineto{\pgfqpoint{1.136036in}{1.369578in}}%
\pgfpathlineto{\pgfqpoint{1.370807in}{1.326904in}}%
\pgfpathlineto{\pgfqpoint{1.371345in}{0.505927in}}%
\pgfpathlineto{\pgfqpoint{1.547466in}{0.505927in}}%
\pgfpathlineto{\pgfqpoint{1.547472in}{1.371382in}}%
\pgfpathlineto{\pgfqpoint{1.548157in}{1.369578in}}%
\pgfpathlineto{\pgfqpoint{1.782928in}{1.326904in}}%
\pgfpathlineto{\pgfqpoint{1.783466in}{0.505927in}}%
\pgfpathlineto{\pgfqpoint{1.959588in}{0.505927in}}%
\pgfpathlineto{\pgfqpoint{1.959593in}{1.371381in}}%
\pgfpathlineto{\pgfqpoint{1.960279in}{1.369578in}}%
\pgfpathlineto{\pgfqpoint{2.195049in}{1.326904in}}%
\pgfpathlineto{\pgfqpoint{2.195587in}{0.505927in}}%
\pgfpathlineto{\pgfqpoint{2.371709in}{0.505927in}}%
\pgfpathlineto{\pgfqpoint{2.371714in}{1.371381in}}%
\pgfpathlineto{\pgfqpoint{2.372400in}{1.369578in}}%
\pgfpathlineto{\pgfqpoint{2.607170in}{1.326904in}}%
\pgfpathlineto{\pgfqpoint{2.607709in}{0.505927in}}%
\pgfpathlineto{\pgfqpoint{2.783830in}{0.505927in}}%
\pgfpathlineto{\pgfqpoint{2.783835in}{1.371381in}}%
\pgfpathlineto{\pgfqpoint{2.784521in}{1.369578in}}%
\pgfpathlineto{\pgfqpoint{3.019281in}{1.326906in}}%
\pgfpathlineto{\pgfqpoint{3.019281in}{1.326906in}}%
\pgfusepath{stroke}%
\end{pgfscope}%
\begin{pgfscope}%
\pgfsetrectcap%
\pgfsetmiterjoin%
\pgfsetlinewidth{0.803000pt}%
\definecolor{currentstroke}{rgb}{0.000000,0.000000,0.000000}%
\pgfsetstrokecolor{currentstroke}%
\pgfsetdash{}{0pt}%
\pgfpathmoveto{\pgfqpoint{0.422918in}{0.462654in}}%
\pgfpathlineto{\pgfqpoint{0.422918in}{1.414654in}}%
\pgfusepath{stroke}%
\end{pgfscope}%
\begin{pgfscope}%
\pgfsetrectcap%
\pgfsetmiterjoin%
\pgfsetlinewidth{0.803000pt}%
\definecolor{currentstroke}{rgb}{0.000000,0.000000,0.000000}%
\pgfsetstrokecolor{currentstroke}%
\pgfsetdash{}{0pt}%
\pgfpathmoveto{\pgfqpoint{3.142917in}{0.462654in}}%
\pgfpathlineto{\pgfqpoint{3.142917in}{1.414654in}}%
\pgfusepath{stroke}%
\end{pgfscope}%
\begin{pgfscope}%
\pgfsetrectcap%
\pgfsetmiterjoin%
\pgfsetlinewidth{0.803000pt}%
\definecolor{currentstroke}{rgb}{0.000000,0.000000,0.000000}%
\pgfsetstrokecolor{currentstroke}%
\pgfsetdash{}{0pt}%
\pgfpathmoveto{\pgfqpoint{0.422918in}{0.462654in}}%
\pgfpathlineto{\pgfqpoint{3.142917in}{0.462654in}}%
\pgfusepath{stroke}%
\end{pgfscope}%
\begin{pgfscope}%
\pgfsetrectcap%
\pgfsetmiterjoin%
\pgfsetlinewidth{0.803000pt}%
\definecolor{currentstroke}{rgb}{0.000000,0.000000,0.000000}%
\pgfsetstrokecolor{currentstroke}%
\pgfsetdash{}{0pt}%
\pgfpathmoveto{\pgfqpoint{0.422918in}{1.414654in}}%
\pgfpathlineto{\pgfqpoint{3.142917in}{1.414654in}}%
\pgfusepath{stroke}%
\end{pgfscope}%
\end{pgfpicture}%
\makeatother%
\endgroup%

	\caption{Corrente do Diodo do circuito Cuk.}
	\label{fig:ex3_1}
\end{figure}

\medskip
\textbf{b)} Calcule o ciclo de trabalho no ponto de operação.

\medskip
\textbf{Resolução}: como o conversor está operando no CCM, pode-se determinar o ciclo de trabalho por~\cite{Eri:20}
\begin{equation*}
	D = \frac{V_o}{V_o + E} = \frac{36}{36 + 48} \approx 0{,}4286
\end{equation*}

\medskip
\textbf{c)} Determine o valor do capacitor intermediário ($C_1$), de modo que a ondulação de tensão sobre ele seja de $0{,}5$ V (pico a pico).

\medskip
\textbf{Resolução}: em conversores que possuem filtro de dois polos (LC) no circuito de saída, o \textit{ripple} da tensão de saída pode ser determinado por~\cite{Eri:20}
\begin{equation}\label{eq:ex3_1}
	\Delta v_{pico} = 2\Delta v = \frac{\Delta i_{L2}T_s}{4C_o},
\end{equation}
em que,
\begin{equation*}
	\Delta i_{L2} = \frac{EDT_s}{2L_1} \approx 16{,}1~\mathrm{mA}.
\end{equation*}

Resolvendo~\eqref{eq:ex3_1}, tem-se que $C_o \approx 1{,}256~\mu$F. Os valores de $\Delta i_{L2}$ e $\Delta v_{pico}$ podem ser verificados na simulação da Figura~\ref{fig:ex3_2}.
\begin{figure}[htb!]
	\centering
	%% Creator: Matplotlib, PGF backend
%%
%% To include the figure in your LaTeX document, write
%%   \input{<filename>.pgf}
%%
%% Make sure the required packages are loaded in your preamble
%%   \usepackage{pgf}
%%
%% Also ensure that all the required font packages are loaded; for instance,
%% the lmodern package is sometimes necessary when using math font.
%%   \usepackage{lmodern}
%%
%% Figures using additional raster images can only be included by \input if
%% they are in the same directory as the main LaTeX file. For loading figures
%% from other directories you can use the `import` package
%%   \usepackage{import}
%%
%% and then include the figures with
%%   \import{<path to file>}{<filename>.pgf}
%%
%% Matplotlib used the following preamble
%%   \def\mathdefault#1{#1}
%%   \everymath=\expandafter{\the\everymath\displaystyle}
%%   \IfFileExists{scrextend.sty}{
%%     \usepackage[fontsize=8.000000pt]{scrextend}
%%   }{
%%     \renewcommand{\normalsize}{\fontsize{8.000000}{9.600000}\selectfont}
%%     \normalsize
%%   }
%%   
%%           \usepackage{amsmath}
%%           \usepackage{amssymb}
%%           \usepackage{mathtools}
%%           \usepackage{dsfont}
%%           \providecommand{\mathdefault}[1]{#1}
%%       
%%   \makeatletter\@ifpackageloaded{underscore}{}{\usepackage[strings]{underscore}}\makeatother
%%
\begingroup%
\makeatletter%
\begin{pgfpicture}%
\pgfpathrectangle{\pgfpointorigin}{\pgfqpoint{3.544619in}{2.322654in}}%
\pgfusepath{use as bounding box, clip}%
\begin{pgfscope}%
\pgfsetbuttcap%
\pgfsetmiterjoin%
\definecolor{currentfill}{rgb}{1.000000,1.000000,1.000000}%
\pgfsetfillcolor{currentfill}%
\pgfsetlinewidth{0.000000pt}%
\definecolor{currentstroke}{rgb}{1.000000,1.000000,1.000000}%
\pgfsetstrokecolor{currentstroke}%
\pgfsetdash{}{0pt}%
\pgfpathmoveto{\pgfqpoint{0.000000in}{0.000000in}}%
\pgfpathlineto{\pgfqpoint{3.544619in}{0.000000in}}%
\pgfpathlineto{\pgfqpoint{3.544619in}{2.322654in}}%
\pgfpathlineto{\pgfqpoint{0.000000in}{2.322654in}}%
\pgfpathlineto{\pgfqpoint{0.000000in}{0.000000in}}%
\pgfpathclose%
\pgfusepath{fill}%
\end{pgfscope}%
\begin{pgfscope}%
\pgfsetbuttcap%
\pgfsetmiterjoin%
\definecolor{currentfill}{rgb}{1.000000,1.000000,1.000000}%
\pgfsetfillcolor{currentfill}%
\pgfsetlinewidth{0.000000pt}%
\definecolor{currentstroke}{rgb}{0.000000,0.000000,0.000000}%
\pgfsetstrokecolor{currentstroke}%
\pgfsetstrokeopacity{0.000000}%
\pgfsetdash{}{0pt}%
\pgfpathmoveto{\pgfqpoint{0.724619in}{1.415315in}}%
\pgfpathlineto{\pgfqpoint{3.444619in}{1.415315in}}%
\pgfpathlineto{\pgfqpoint{3.444619in}{2.222654in}}%
\pgfpathlineto{\pgfqpoint{0.724619in}{2.222654in}}%
\pgfpathlineto{\pgfqpoint{0.724619in}{1.415315in}}%
\pgfpathclose%
\pgfusepath{fill}%
\end{pgfscope}%
\begin{pgfscope}%
\pgfpathrectangle{\pgfqpoint{0.724619in}{1.415315in}}{\pgfqpoint{2.720000in}{0.807339in}}%
\pgfusepath{clip}%
\pgfsetbuttcap%
\pgfsetroundjoin%
\pgfsetlinewidth{0.501875pt}%
\definecolor{currentstroke}{rgb}{0.690196,0.690196,0.690196}%
\pgfsetstrokecolor{currentstroke}%
\pgfsetstrokeopacity{0.250000}%
\pgfsetdash{{1.850000pt}{0.800000pt}}{0.000000pt}%
\pgfpathmoveto{\pgfqpoint{1.079044in}{1.415315in}}%
\pgfpathlineto{\pgfqpoint{1.079044in}{2.222654in}}%
\pgfusepath{stroke}%
\end{pgfscope}%
\begin{pgfscope}%
\pgfsetbuttcap%
\pgfsetroundjoin%
\definecolor{currentfill}{rgb}{0.000000,0.000000,0.000000}%
\pgfsetfillcolor{currentfill}%
\pgfsetlinewidth{0.803000pt}%
\definecolor{currentstroke}{rgb}{0.000000,0.000000,0.000000}%
\pgfsetstrokecolor{currentstroke}%
\pgfsetdash{}{0pt}%
\pgfsys@defobject{currentmarker}{\pgfqpoint{0.000000in}{-0.048611in}}{\pgfqpoint{0.000000in}{0.000000in}}{%
\pgfpathmoveto{\pgfqpoint{0.000000in}{0.000000in}}%
\pgfpathlineto{\pgfqpoint{0.000000in}{-0.048611in}}%
\pgfusepath{stroke,fill}%
}%
\begin{pgfscope}%
\pgfsys@transformshift{1.079044in}{1.415315in}%
\pgfsys@useobject{currentmarker}{}%
\end{pgfscope}%
\end{pgfscope}%
\begin{pgfscope}%
\pgfpathrectangle{\pgfqpoint{0.724619in}{1.415315in}}{\pgfqpoint{2.720000in}{0.807339in}}%
\pgfusepath{clip}%
\pgfsetbuttcap%
\pgfsetroundjoin%
\pgfsetlinewidth{0.501875pt}%
\definecolor{currentstroke}{rgb}{0.690196,0.690196,0.690196}%
\pgfsetstrokecolor{currentstroke}%
\pgfsetstrokeopacity{0.250000}%
\pgfsetdash{{1.850000pt}{0.800000pt}}{0.000000pt}%
\pgfpathmoveto{\pgfqpoint{1.606559in}{1.415315in}}%
\pgfpathlineto{\pgfqpoint{1.606559in}{2.222654in}}%
\pgfusepath{stroke}%
\end{pgfscope}%
\begin{pgfscope}%
\pgfsetbuttcap%
\pgfsetroundjoin%
\definecolor{currentfill}{rgb}{0.000000,0.000000,0.000000}%
\pgfsetfillcolor{currentfill}%
\pgfsetlinewidth{0.803000pt}%
\definecolor{currentstroke}{rgb}{0.000000,0.000000,0.000000}%
\pgfsetstrokecolor{currentstroke}%
\pgfsetdash{}{0pt}%
\pgfsys@defobject{currentmarker}{\pgfqpoint{0.000000in}{-0.048611in}}{\pgfqpoint{0.000000in}{0.000000in}}{%
\pgfpathmoveto{\pgfqpoint{0.000000in}{0.000000in}}%
\pgfpathlineto{\pgfqpoint{0.000000in}{-0.048611in}}%
\pgfusepath{stroke,fill}%
}%
\begin{pgfscope}%
\pgfsys@transformshift{1.606559in}{1.415315in}%
\pgfsys@useobject{currentmarker}{}%
\end{pgfscope}%
\end{pgfscope}%
\begin{pgfscope}%
\pgfpathrectangle{\pgfqpoint{0.724619in}{1.415315in}}{\pgfqpoint{2.720000in}{0.807339in}}%
\pgfusepath{clip}%
\pgfsetbuttcap%
\pgfsetroundjoin%
\pgfsetlinewidth{0.501875pt}%
\definecolor{currentstroke}{rgb}{0.690196,0.690196,0.690196}%
\pgfsetstrokecolor{currentstroke}%
\pgfsetstrokeopacity{0.250000}%
\pgfsetdash{{1.850000pt}{0.800000pt}}{0.000000pt}%
\pgfpathmoveto{\pgfqpoint{2.134074in}{1.415315in}}%
\pgfpathlineto{\pgfqpoint{2.134074in}{2.222654in}}%
\pgfusepath{stroke}%
\end{pgfscope}%
\begin{pgfscope}%
\pgfsetbuttcap%
\pgfsetroundjoin%
\definecolor{currentfill}{rgb}{0.000000,0.000000,0.000000}%
\pgfsetfillcolor{currentfill}%
\pgfsetlinewidth{0.803000pt}%
\definecolor{currentstroke}{rgb}{0.000000,0.000000,0.000000}%
\pgfsetstrokecolor{currentstroke}%
\pgfsetdash{}{0pt}%
\pgfsys@defobject{currentmarker}{\pgfqpoint{0.000000in}{-0.048611in}}{\pgfqpoint{0.000000in}{0.000000in}}{%
\pgfpathmoveto{\pgfqpoint{0.000000in}{0.000000in}}%
\pgfpathlineto{\pgfqpoint{0.000000in}{-0.048611in}}%
\pgfusepath{stroke,fill}%
}%
\begin{pgfscope}%
\pgfsys@transformshift{2.134074in}{1.415315in}%
\pgfsys@useobject{currentmarker}{}%
\end{pgfscope}%
\end{pgfscope}%
\begin{pgfscope}%
\pgfpathrectangle{\pgfqpoint{0.724619in}{1.415315in}}{\pgfqpoint{2.720000in}{0.807339in}}%
\pgfusepath{clip}%
\pgfsetbuttcap%
\pgfsetroundjoin%
\pgfsetlinewidth{0.501875pt}%
\definecolor{currentstroke}{rgb}{0.690196,0.690196,0.690196}%
\pgfsetstrokecolor{currentstroke}%
\pgfsetstrokeopacity{0.250000}%
\pgfsetdash{{1.850000pt}{0.800000pt}}{0.000000pt}%
\pgfpathmoveto{\pgfqpoint{2.661589in}{1.415315in}}%
\pgfpathlineto{\pgfqpoint{2.661589in}{2.222654in}}%
\pgfusepath{stroke}%
\end{pgfscope}%
\begin{pgfscope}%
\pgfsetbuttcap%
\pgfsetroundjoin%
\definecolor{currentfill}{rgb}{0.000000,0.000000,0.000000}%
\pgfsetfillcolor{currentfill}%
\pgfsetlinewidth{0.803000pt}%
\definecolor{currentstroke}{rgb}{0.000000,0.000000,0.000000}%
\pgfsetstrokecolor{currentstroke}%
\pgfsetdash{}{0pt}%
\pgfsys@defobject{currentmarker}{\pgfqpoint{0.000000in}{-0.048611in}}{\pgfqpoint{0.000000in}{0.000000in}}{%
\pgfpathmoveto{\pgfqpoint{0.000000in}{0.000000in}}%
\pgfpathlineto{\pgfqpoint{0.000000in}{-0.048611in}}%
\pgfusepath{stroke,fill}%
}%
\begin{pgfscope}%
\pgfsys@transformshift{2.661589in}{1.415315in}%
\pgfsys@useobject{currentmarker}{}%
\end{pgfscope}%
\end{pgfscope}%
\begin{pgfscope}%
\pgfpathrectangle{\pgfqpoint{0.724619in}{1.415315in}}{\pgfqpoint{2.720000in}{0.807339in}}%
\pgfusepath{clip}%
\pgfsetbuttcap%
\pgfsetroundjoin%
\pgfsetlinewidth{0.501875pt}%
\definecolor{currentstroke}{rgb}{0.690196,0.690196,0.690196}%
\pgfsetstrokecolor{currentstroke}%
\pgfsetstrokeopacity{0.250000}%
\pgfsetdash{{1.850000pt}{0.800000pt}}{0.000000pt}%
\pgfpathmoveto{\pgfqpoint{3.189104in}{1.415315in}}%
\pgfpathlineto{\pgfqpoint{3.189104in}{2.222654in}}%
\pgfusepath{stroke}%
\end{pgfscope}%
\begin{pgfscope}%
\pgfsetbuttcap%
\pgfsetroundjoin%
\definecolor{currentfill}{rgb}{0.000000,0.000000,0.000000}%
\pgfsetfillcolor{currentfill}%
\pgfsetlinewidth{0.803000pt}%
\definecolor{currentstroke}{rgb}{0.000000,0.000000,0.000000}%
\pgfsetstrokecolor{currentstroke}%
\pgfsetdash{}{0pt}%
\pgfsys@defobject{currentmarker}{\pgfqpoint{0.000000in}{-0.048611in}}{\pgfqpoint{0.000000in}{0.000000in}}{%
\pgfpathmoveto{\pgfqpoint{0.000000in}{0.000000in}}%
\pgfpathlineto{\pgfqpoint{0.000000in}{-0.048611in}}%
\pgfusepath{stroke,fill}%
}%
\begin{pgfscope}%
\pgfsys@transformshift{3.189104in}{1.415315in}%
\pgfsys@useobject{currentmarker}{}%
\end{pgfscope}%
\end{pgfscope}%
\begin{pgfscope}%
\pgfpathrectangle{\pgfqpoint{0.724619in}{1.415315in}}{\pgfqpoint{2.720000in}{0.807339in}}%
\pgfusepath{clip}%
\pgfsetbuttcap%
\pgfsetroundjoin%
\pgfsetlinewidth{0.501875pt}%
\definecolor{currentstroke}{rgb}{0.690196,0.690196,0.690196}%
\pgfsetstrokecolor{currentstroke}%
\pgfsetstrokeopacity{0.250000}%
\pgfsetdash{{1.850000pt}{0.800000pt}}{0.000000pt}%
\pgfpathmoveto{\pgfqpoint{0.724619in}{1.481100in}}%
\pgfpathlineto{\pgfqpoint{3.444619in}{1.481100in}}%
\pgfusepath{stroke}%
\end{pgfscope}%
\begin{pgfscope}%
\pgfsetbuttcap%
\pgfsetroundjoin%
\definecolor{currentfill}{rgb}{0.000000,0.000000,0.000000}%
\pgfsetfillcolor{currentfill}%
\pgfsetlinewidth{0.803000pt}%
\definecolor{currentstroke}{rgb}{0.000000,0.000000,0.000000}%
\pgfsetstrokecolor{currentstroke}%
\pgfsetdash{}{0pt}%
\pgfsys@defobject{currentmarker}{\pgfqpoint{-0.048611in}{0.000000in}}{\pgfqpoint{-0.000000in}{0.000000in}}{%
\pgfpathmoveto{\pgfqpoint{-0.000000in}{0.000000in}}%
\pgfpathlineto{\pgfqpoint{-0.048611in}{0.000000in}}%
\pgfusepath{stroke,fill}%
}%
\begin{pgfscope}%
\pgfsys@transformshift{0.724619in}{1.481100in}%
\pgfsys@useobject{currentmarker}{}%
\end{pgfscope}%
\end{pgfscope}%
\begin{pgfscope}%
\definecolor{textcolor}{rgb}{0.000000,0.000000,0.000000}%
\pgfsetstrokecolor{textcolor}%
\pgfsetfillcolor{textcolor}%
\pgftext[x=0.266667in, y=1.442520in, left, base]{\color{textcolor}{\rmfamily\fontsize{8.000000}{9.600000}\selectfont\catcode`\^=\active\def^{\ifmmode\sp\else\^{}\fi}\catcode`\%=\active\def%{\%}$\mathdefault{\ensuremath{-}36.00}$}}%
\end{pgfscope}%
\begin{pgfscope}%
\pgfpathrectangle{\pgfqpoint{0.724619in}{1.415315in}}{\pgfqpoint{2.720000in}{0.807339in}}%
\pgfusepath{clip}%
\pgfsetbuttcap%
\pgfsetroundjoin%
\pgfsetlinewidth{0.501875pt}%
\definecolor{currentstroke}{rgb}{0.690196,0.690196,0.690196}%
\pgfsetstrokecolor{currentstroke}%
\pgfsetstrokeopacity{0.250000}%
\pgfsetdash{{1.850000pt}{0.800000pt}}{0.000000pt}%
\pgfpathmoveto{\pgfqpoint{0.724619in}{1.704628in}}%
\pgfpathlineto{\pgfqpoint{3.444619in}{1.704628in}}%
\pgfusepath{stroke}%
\end{pgfscope}%
\begin{pgfscope}%
\pgfsetbuttcap%
\pgfsetroundjoin%
\definecolor{currentfill}{rgb}{0.000000,0.000000,0.000000}%
\pgfsetfillcolor{currentfill}%
\pgfsetlinewidth{0.803000pt}%
\definecolor{currentstroke}{rgb}{0.000000,0.000000,0.000000}%
\pgfsetstrokecolor{currentstroke}%
\pgfsetdash{}{0pt}%
\pgfsys@defobject{currentmarker}{\pgfqpoint{-0.048611in}{0.000000in}}{\pgfqpoint{-0.000000in}{0.000000in}}{%
\pgfpathmoveto{\pgfqpoint{-0.000000in}{0.000000in}}%
\pgfpathlineto{\pgfqpoint{-0.048611in}{0.000000in}}%
\pgfusepath{stroke,fill}%
}%
\begin{pgfscope}%
\pgfsys@transformshift{0.724619in}{1.704628in}%
\pgfsys@useobject{currentmarker}{}%
\end{pgfscope}%
\end{pgfscope}%
\begin{pgfscope}%
\definecolor{textcolor}{rgb}{0.000000,0.000000,0.000000}%
\pgfsetstrokecolor{textcolor}%
\pgfsetfillcolor{textcolor}%
\pgftext[x=0.266667in, y=1.666048in, left, base]{\color{textcolor}{\rmfamily\fontsize{8.000000}{9.600000}\selectfont\catcode`\^=\active\def^{\ifmmode\sp\else\^{}\fi}\catcode`\%=\active\def%{\%}$\mathdefault{\ensuremath{-}35.85}$}}%
\end{pgfscope}%
\begin{pgfscope}%
\pgfpathrectangle{\pgfqpoint{0.724619in}{1.415315in}}{\pgfqpoint{2.720000in}{0.807339in}}%
\pgfusepath{clip}%
\pgfsetbuttcap%
\pgfsetroundjoin%
\pgfsetlinewidth{0.501875pt}%
\definecolor{currentstroke}{rgb}{0.690196,0.690196,0.690196}%
\pgfsetstrokecolor{currentstroke}%
\pgfsetstrokeopacity{0.250000}%
\pgfsetdash{{1.850000pt}{0.800000pt}}{0.000000pt}%
\pgfpathmoveto{\pgfqpoint{0.724619in}{1.928156in}}%
\pgfpathlineto{\pgfqpoint{3.444619in}{1.928156in}}%
\pgfusepath{stroke}%
\end{pgfscope}%
\begin{pgfscope}%
\pgfsetbuttcap%
\pgfsetroundjoin%
\definecolor{currentfill}{rgb}{0.000000,0.000000,0.000000}%
\pgfsetfillcolor{currentfill}%
\pgfsetlinewidth{0.803000pt}%
\definecolor{currentstroke}{rgb}{0.000000,0.000000,0.000000}%
\pgfsetstrokecolor{currentstroke}%
\pgfsetdash{}{0pt}%
\pgfsys@defobject{currentmarker}{\pgfqpoint{-0.048611in}{0.000000in}}{\pgfqpoint{-0.000000in}{0.000000in}}{%
\pgfpathmoveto{\pgfqpoint{-0.000000in}{0.000000in}}%
\pgfpathlineto{\pgfqpoint{-0.048611in}{0.000000in}}%
\pgfusepath{stroke,fill}%
}%
\begin{pgfscope}%
\pgfsys@transformshift{0.724619in}{1.928156in}%
\pgfsys@useobject{currentmarker}{}%
\end{pgfscope}%
\end{pgfscope}%
\begin{pgfscope}%
\definecolor{textcolor}{rgb}{0.000000,0.000000,0.000000}%
\pgfsetstrokecolor{textcolor}%
\pgfsetfillcolor{textcolor}%
\pgftext[x=0.266667in, y=1.889575in, left, base]{\color{textcolor}{\rmfamily\fontsize{8.000000}{9.600000}\selectfont\catcode`\^=\active\def^{\ifmmode\sp\else\^{}\fi}\catcode`\%=\active\def%{\%}$\mathdefault{\ensuremath{-}35.70}$}}%
\end{pgfscope}%
\begin{pgfscope}%
\pgfpathrectangle{\pgfqpoint{0.724619in}{1.415315in}}{\pgfqpoint{2.720000in}{0.807339in}}%
\pgfusepath{clip}%
\pgfsetbuttcap%
\pgfsetroundjoin%
\pgfsetlinewidth{0.501875pt}%
\definecolor{currentstroke}{rgb}{0.690196,0.690196,0.690196}%
\pgfsetstrokecolor{currentstroke}%
\pgfsetstrokeopacity{0.250000}%
\pgfsetdash{{1.850000pt}{0.800000pt}}{0.000000pt}%
\pgfpathmoveto{\pgfqpoint{0.724619in}{2.151683in}}%
\pgfpathlineto{\pgfqpoint{3.444619in}{2.151683in}}%
\pgfusepath{stroke}%
\end{pgfscope}%
\begin{pgfscope}%
\pgfsetbuttcap%
\pgfsetroundjoin%
\definecolor{currentfill}{rgb}{0.000000,0.000000,0.000000}%
\pgfsetfillcolor{currentfill}%
\pgfsetlinewidth{0.803000pt}%
\definecolor{currentstroke}{rgb}{0.000000,0.000000,0.000000}%
\pgfsetstrokecolor{currentstroke}%
\pgfsetdash{}{0pt}%
\pgfsys@defobject{currentmarker}{\pgfqpoint{-0.048611in}{0.000000in}}{\pgfqpoint{-0.000000in}{0.000000in}}{%
\pgfpathmoveto{\pgfqpoint{-0.000000in}{0.000000in}}%
\pgfpathlineto{\pgfqpoint{-0.048611in}{0.000000in}}%
\pgfusepath{stroke,fill}%
}%
\begin{pgfscope}%
\pgfsys@transformshift{0.724619in}{2.151683in}%
\pgfsys@useobject{currentmarker}{}%
\end{pgfscope}%
\end{pgfscope}%
\begin{pgfscope}%
\definecolor{textcolor}{rgb}{0.000000,0.000000,0.000000}%
\pgfsetstrokecolor{textcolor}%
\pgfsetfillcolor{textcolor}%
\pgftext[x=0.266667in, y=2.113103in, left, base]{\color{textcolor}{\rmfamily\fontsize{8.000000}{9.600000}\selectfont\catcode`\^=\active\def^{\ifmmode\sp\else\^{}\fi}\catcode`\%=\active\def%{\%}$\mathdefault{\ensuremath{-}35.55}$}}%
\end{pgfscope}%
\begin{pgfscope}%
\definecolor{textcolor}{rgb}{0.000000,0.000000,0.000000}%
\pgfsetstrokecolor{textcolor}%
\pgfsetfillcolor{textcolor}%
\pgftext[x=0.211111in,y=1.818985in,,bottom,rotate=90.000000]{\color{textcolor}{\rmfamily\fontsize{8.000000}{9.600000}\selectfont\catcode`\^=\active\def^{\ifmmode\sp\else\^{}\fi}\catcode`\%=\active\def%{\%}$v_o$ (V)}}%
\end{pgfscope}%
\begin{pgfscope}%
\pgfpathrectangle{\pgfqpoint{0.724619in}{1.415315in}}{\pgfqpoint{2.720000in}{0.807339in}}%
\pgfusepath{clip}%
\pgfsetrectcap%
\pgfsetroundjoin%
\pgfsetlinewidth{1.505625pt}%
\definecolor{currentstroke}{rgb}{0.000000,0.447059,0.698039}%
\pgfsetstrokecolor{currentstroke}%
\pgfsetdash{}{0pt}%
\pgfpathmoveto{\pgfqpoint{0.848256in}{1.950688in}}%
\pgfpathlineto{\pgfqpoint{0.856588in}{2.003430in}}%
\pgfpathlineto{\pgfqpoint{0.864830in}{2.048626in}}%
\pgfpathlineto{\pgfqpoint{0.873073in}{2.087074in}}%
\pgfpathlineto{\pgfqpoint{0.881315in}{2.118955in}}%
\pgfpathlineto{\pgfqpoint{0.889558in}{2.144447in}}%
\pgfpathlineto{\pgfqpoint{0.897800in}{2.163724in}}%
\pgfpathlineto{\pgfqpoint{0.906042in}{2.176956in}}%
\pgfpathlineto{\pgfqpoint{0.910165in}{2.181357in}}%
\pgfpathlineto{\pgfqpoint{0.914285in}{2.184308in}}%
\pgfpathlineto{\pgfqpoint{0.918407in}{2.185829in}}%
\pgfpathlineto{\pgfqpoint{0.922527in}{2.185940in}}%
\pgfpathlineto{\pgfqpoint{0.926650in}{2.184661in}}%
\pgfpathlineto{\pgfqpoint{0.930770in}{2.182011in}}%
\pgfpathlineto{\pgfqpoint{0.934892in}{2.178008in}}%
\pgfpathlineto{\pgfqpoint{0.939012in}{2.172672in}}%
\pgfpathlineto{\pgfqpoint{0.947255in}{2.158075in}}%
\pgfpathlineto{\pgfqpoint{0.955497in}{2.138365in}}%
\pgfpathlineto{\pgfqpoint{0.963739in}{2.113684in}}%
\pgfpathlineto{\pgfqpoint{0.971982in}{2.084172in}}%
\pgfpathlineto{\pgfqpoint{0.980224in}{2.049963in}}%
\pgfpathlineto{\pgfqpoint{0.992589in}{1.990132in}}%
\pgfpathlineto{\pgfqpoint{1.004952in}{1.920463in}}%
\pgfpathlineto{\pgfqpoint{1.017317in}{1.841373in}}%
\pgfpathlineto{\pgfqpoint{1.038586in}{1.697547in}}%
\pgfpathlineto{\pgfqpoint{1.050951in}{1.628711in}}%
\pgfpathlineto{\pgfqpoint{1.059193in}{1.589529in}}%
\pgfpathlineto{\pgfqpoint{1.067436in}{1.555548in}}%
\pgfpathlineto{\pgfqpoint{1.075678in}{1.526628in}}%
\pgfpathlineto{\pgfqpoint{1.083921in}{1.502633in}}%
\pgfpathlineto{\pgfqpoint{1.092163in}{1.483430in}}%
\pgfpathlineto{\pgfqpoint{1.100405in}{1.468892in}}%
\pgfpathlineto{\pgfqpoint{1.108648in}{1.458891in}}%
\pgfpathlineto{\pgfqpoint{1.112768in}{1.455553in}}%
\pgfpathlineto{\pgfqpoint{1.116890in}{1.453304in}}%
\pgfpathlineto{\pgfqpoint{1.121010in}{1.452129in}}%
\pgfpathlineto{\pgfqpoint{1.125133in}{1.452012in}}%
\pgfpathlineto{\pgfqpoint{1.129253in}{1.452940in}}%
\pgfpathlineto{\pgfqpoint{1.133375in}{1.454898in}}%
\pgfpathlineto{\pgfqpoint{1.137495in}{1.457872in}}%
\pgfpathlineto{\pgfqpoint{1.145737in}{1.466811in}}%
\pgfpathlineto{\pgfqpoint{1.153980in}{1.479648in}}%
\pgfpathlineto{\pgfqpoint{1.162222in}{1.496274in}}%
\pgfpathlineto{\pgfqpoint{1.170465in}{1.516587in}}%
\pgfpathlineto{\pgfqpoint{1.178707in}{1.540483in}}%
\pgfpathlineto{\pgfqpoint{1.191072in}{1.582830in}}%
\pgfpathlineto{\pgfqpoint{1.203434in}{1.632690in}}%
\pgfpathlineto{\pgfqpoint{1.215799in}{1.689750in}}%
\pgfpathlineto{\pgfqpoint{1.228162in}{1.753708in}}%
\pgfpathlineto{\pgfqpoint{1.244647in}{1.849214in}}%
\pgfpathlineto{\pgfqpoint{1.262495in}{1.964782in}}%
\pgfpathlineto{\pgfqpoint{1.272832in}{2.026887in}}%
\pgfpathlineto{\pgfqpoint{1.281074in}{2.068686in}}%
\pgfpathlineto{\pgfqpoint{1.289317in}{2.103827in}}%
\pgfpathlineto{\pgfqpoint{1.297559in}{2.132492in}}%
\pgfpathlineto{\pgfqpoint{1.305801in}{2.154855in}}%
\pgfpathlineto{\pgfqpoint{1.314044in}{2.171089in}}%
\pgfpathlineto{\pgfqpoint{1.318164in}{2.176960in}}%
\pgfpathlineto{\pgfqpoint{1.322286in}{2.181360in}}%
\pgfpathlineto{\pgfqpoint{1.326406in}{2.184311in}}%
\pgfpathlineto{\pgfqpoint{1.330529in}{2.185832in}}%
\pgfpathlineto{\pgfqpoint{1.334649in}{2.185943in}}%
\pgfpathlineto{\pgfqpoint{1.338771in}{2.184664in}}%
\pgfpathlineto{\pgfqpoint{1.342891in}{2.182014in}}%
\pgfpathlineto{\pgfqpoint{1.347013in}{2.178012in}}%
\pgfpathlineto{\pgfqpoint{1.351133in}{2.172676in}}%
\pgfpathlineto{\pgfqpoint{1.359376in}{2.158079in}}%
\pgfpathlineto{\pgfqpoint{1.367618in}{2.138369in}}%
\pgfpathlineto{\pgfqpoint{1.375861in}{2.113688in}}%
\pgfpathlineto{\pgfqpoint{1.384103in}{2.084175in}}%
\pgfpathlineto{\pgfqpoint{1.392345in}{2.049966in}}%
\pgfpathlineto{\pgfqpoint{1.404710in}{1.990135in}}%
\pgfpathlineto{\pgfqpoint{1.417073in}{1.920467in}}%
\pgfpathlineto{\pgfqpoint{1.429438in}{1.841376in}}%
\pgfpathlineto{\pgfqpoint{1.450707in}{1.697551in}}%
\pgfpathlineto{\pgfqpoint{1.463072in}{1.628715in}}%
\pgfpathlineto{\pgfqpoint{1.471315in}{1.589533in}}%
\pgfpathlineto{\pgfqpoint{1.479557in}{1.555551in}}%
\pgfpathlineto{\pgfqpoint{1.487799in}{1.526631in}}%
\pgfpathlineto{\pgfqpoint{1.496042in}{1.502636in}}%
\pgfpathlineto{\pgfqpoint{1.504284in}{1.483434in}}%
\pgfpathlineto{\pgfqpoint{1.512527in}{1.468895in}}%
\pgfpathlineto{\pgfqpoint{1.520769in}{1.458894in}}%
\pgfpathlineto{\pgfqpoint{1.524889in}{1.455556in}}%
\pgfpathlineto{\pgfqpoint{1.529011in}{1.453308in}}%
\pgfpathlineto{\pgfqpoint{1.533131in}{1.452132in}}%
\pgfpathlineto{\pgfqpoint{1.537254in}{1.452016in}}%
\pgfpathlineto{\pgfqpoint{1.541374in}{1.452944in}}%
\pgfpathlineto{\pgfqpoint{1.545496in}{1.454902in}}%
\pgfpathlineto{\pgfqpoint{1.549616in}{1.457875in}}%
\pgfpathlineto{\pgfqpoint{1.557859in}{1.466814in}}%
\pgfpathlineto{\pgfqpoint{1.566101in}{1.479651in}}%
\pgfpathlineto{\pgfqpoint{1.574343in}{1.496278in}}%
\pgfpathlineto{\pgfqpoint{1.582586in}{1.516591in}}%
\pgfpathlineto{\pgfqpoint{1.590828in}{1.540487in}}%
\pgfpathlineto{\pgfqpoint{1.603193in}{1.582833in}}%
\pgfpathlineto{\pgfqpoint{1.615556in}{1.632693in}}%
\pgfpathlineto{\pgfqpoint{1.627921in}{1.689754in}}%
\pgfpathlineto{\pgfqpoint{1.640283in}{1.753712in}}%
\pgfpathlineto{\pgfqpoint{1.656768in}{1.849218in}}%
\pgfpathlineto{\pgfqpoint{1.674616in}{1.964786in}}%
\pgfpathlineto{\pgfqpoint{1.684953in}{2.026890in}}%
\pgfpathlineto{\pgfqpoint{1.693195in}{2.068689in}}%
\pgfpathlineto{\pgfqpoint{1.701438in}{2.103831in}}%
\pgfpathlineto{\pgfqpoint{1.709680in}{2.132495in}}%
\pgfpathlineto{\pgfqpoint{1.717923in}{2.154859in}}%
\pgfpathlineto{\pgfqpoint{1.726165in}{2.171092in}}%
\pgfpathlineto{\pgfqpoint{1.730285in}{2.176963in}}%
\pgfpathlineto{\pgfqpoint{1.734407in}{2.181364in}}%
\pgfpathlineto{\pgfqpoint{1.738527in}{2.184315in}}%
\pgfpathlineto{\pgfqpoint{1.742650in}{2.185836in}}%
\pgfpathlineto{\pgfqpoint{1.746770in}{2.185947in}}%
\pgfpathlineto{\pgfqpoint{1.750892in}{2.184668in}}%
\pgfpathlineto{\pgfqpoint{1.755012in}{2.182017in}}%
\pgfpathlineto{\pgfqpoint{1.759135in}{2.178015in}}%
\pgfpathlineto{\pgfqpoint{1.763255in}{2.172679in}}%
\pgfpathlineto{\pgfqpoint{1.771497in}{2.158082in}}%
\pgfpathlineto{\pgfqpoint{1.779739in}{2.138372in}}%
\pgfpathlineto{\pgfqpoint{1.787982in}{2.113691in}}%
\pgfpathlineto{\pgfqpoint{1.796224in}{2.084179in}}%
\pgfpathlineto{\pgfqpoint{1.804467in}{2.049970in}}%
\pgfpathlineto{\pgfqpoint{1.816832in}{1.990139in}}%
\pgfpathlineto{\pgfqpoint{1.829194in}{1.920470in}}%
\pgfpathlineto{\pgfqpoint{1.841559in}{1.841380in}}%
\pgfpathlineto{\pgfqpoint{1.862828in}{1.697554in}}%
\pgfpathlineto{\pgfqpoint{1.875193in}{1.628718in}}%
\pgfpathlineto{\pgfqpoint{1.883436in}{1.589536in}}%
\pgfpathlineto{\pgfqpoint{1.891678in}{1.555555in}}%
\pgfpathlineto{\pgfqpoint{1.899921in}{1.526635in}}%
\pgfpathlineto{\pgfqpoint{1.908163in}{1.502639in}}%
\pgfpathlineto{\pgfqpoint{1.916405in}{1.483437in}}%
\pgfpathlineto{\pgfqpoint{1.924648in}{1.468899in}}%
\pgfpathlineto{\pgfqpoint{1.932890in}{1.458897in}}%
\pgfpathlineto{\pgfqpoint{1.937010in}{1.455560in}}%
\pgfpathlineto{\pgfqpoint{1.941133in}{1.453311in}}%
\pgfpathlineto{\pgfqpoint{1.945253in}{1.452136in}}%
\pgfpathlineto{\pgfqpoint{1.949375in}{1.452019in}}%
\pgfpathlineto{\pgfqpoint{1.953495in}{1.452947in}}%
\pgfpathlineto{\pgfqpoint{1.957618in}{1.454905in}}%
\pgfpathlineto{\pgfqpoint{1.961737in}{1.457879in}}%
\pgfpathlineto{\pgfqpoint{1.969980in}{1.466818in}}%
\pgfpathlineto{\pgfqpoint{1.978222in}{1.479655in}}%
\pgfpathlineto{\pgfqpoint{1.986465in}{1.496281in}}%
\pgfpathlineto{\pgfqpoint{1.994707in}{1.516594in}}%
\pgfpathlineto{\pgfqpoint{2.002950in}{1.540490in}}%
\pgfpathlineto{\pgfqpoint{2.015315in}{1.582836in}}%
\pgfpathlineto{\pgfqpoint{2.027677in}{1.632697in}}%
\pgfpathlineto{\pgfqpoint{2.040042in}{1.689757in}}%
\pgfpathlineto{\pgfqpoint{2.052404in}{1.753715in}}%
\pgfpathlineto{\pgfqpoint{2.068889in}{1.849221in}}%
\pgfpathlineto{\pgfqpoint{2.086737in}{1.964789in}}%
\pgfpathlineto{\pgfqpoint{2.097074in}{2.026894in}}%
\pgfpathlineto{\pgfqpoint{2.105317in}{2.068693in}}%
\pgfpathlineto{\pgfqpoint{2.113559in}{2.103834in}}%
\pgfpathlineto{\pgfqpoint{2.121801in}{2.132499in}}%
\pgfpathlineto{\pgfqpoint{2.130044in}{2.154862in}}%
\pgfpathlineto{\pgfqpoint{2.138286in}{2.171096in}}%
\pgfpathlineto{\pgfqpoint{2.142406in}{2.176966in}}%
\pgfpathlineto{\pgfqpoint{2.146529in}{2.181367in}}%
\pgfpathlineto{\pgfqpoint{2.150649in}{2.184318in}}%
\pgfpathlineto{\pgfqpoint{2.154771in}{2.185839in}}%
\pgfpathlineto{\pgfqpoint{2.158891in}{2.185950in}}%
\pgfpathlineto{\pgfqpoint{2.163013in}{2.184671in}}%
\pgfpathlineto{\pgfqpoint{2.167133in}{2.182021in}}%
\pgfpathlineto{\pgfqpoint{2.171256in}{2.178018in}}%
\pgfpathlineto{\pgfqpoint{2.175376in}{2.172683in}}%
\pgfpathlineto{\pgfqpoint{2.183618in}{2.158086in}}%
\pgfpathlineto{\pgfqpoint{2.191861in}{2.138375in}}%
\pgfpathlineto{\pgfqpoint{2.200103in}{2.113695in}}%
\pgfpathlineto{\pgfqpoint{2.208345in}{2.084182in}}%
\pgfpathlineto{\pgfqpoint{2.216588in}{2.049973in}}%
\pgfpathlineto{\pgfqpoint{2.228953in}{1.990142in}}%
\pgfpathlineto{\pgfqpoint{2.241315in}{1.920473in}}%
\pgfpathlineto{\pgfqpoint{2.253680in}{1.841383in}}%
\pgfpathlineto{\pgfqpoint{2.274950in}{1.697558in}}%
\pgfpathlineto{\pgfqpoint{2.287315in}{1.628721in}}%
\pgfpathlineto{\pgfqpoint{2.295557in}{1.589540in}}%
\pgfpathlineto{\pgfqpoint{2.303799in}{1.555558in}}%
\pgfpathlineto{\pgfqpoint{2.312042in}{1.526638in}}%
\pgfpathlineto{\pgfqpoint{2.320284in}{1.502643in}}%
\pgfpathlineto{\pgfqpoint{2.328527in}{1.483441in}}%
\pgfpathlineto{\pgfqpoint{2.336769in}{1.468902in}}%
\pgfpathlineto{\pgfqpoint{2.345011in}{1.458901in}}%
\pgfpathlineto{\pgfqpoint{2.349131in}{1.455564in}}%
\pgfpathlineto{\pgfqpoint{2.353254in}{1.453315in}}%
\pgfpathlineto{\pgfqpoint{2.357374in}{1.452139in}}%
\pgfpathlineto{\pgfqpoint{2.361496in}{1.452023in}}%
\pgfpathlineto{\pgfqpoint{2.365616in}{1.452951in}}%
\pgfpathlineto{\pgfqpoint{2.369739in}{1.454908in}}%
\pgfpathlineto{\pgfqpoint{2.373859in}{1.457882in}}%
\pgfpathlineto{\pgfqpoint{2.382101in}{1.466821in}}%
\pgfpathlineto{\pgfqpoint{2.390343in}{1.479658in}}%
\pgfpathlineto{\pgfqpoint{2.398586in}{1.496285in}}%
\pgfpathlineto{\pgfqpoint{2.406828in}{1.516597in}}%
\pgfpathlineto{\pgfqpoint{2.415071in}{1.540494in}}%
\pgfpathlineto{\pgfqpoint{2.427436in}{1.582840in}}%
\pgfpathlineto{\pgfqpoint{2.439798in}{1.632700in}}%
\pgfpathlineto{\pgfqpoint{2.452163in}{1.689760in}}%
\pgfpathlineto{\pgfqpoint{2.464525in}{1.753719in}}%
\pgfpathlineto{\pgfqpoint{2.481010in}{1.849224in}}%
\pgfpathlineto{\pgfqpoint{2.498859in}{1.964793in}}%
\pgfpathlineto{\pgfqpoint{2.509195in}{2.026897in}}%
\pgfpathlineto{\pgfqpoint{2.517438in}{2.068696in}}%
\pgfpathlineto{\pgfqpoint{2.525680in}{2.103837in}}%
\pgfpathlineto{\pgfqpoint{2.533923in}{2.132502in}}%
\pgfpathlineto{\pgfqpoint{2.542165in}{2.154866in}}%
\pgfpathlineto{\pgfqpoint{2.550407in}{2.171099in}}%
\pgfpathlineto{\pgfqpoint{2.554527in}{2.176970in}}%
\pgfpathlineto{\pgfqpoint{2.558650in}{2.181371in}}%
\pgfpathlineto{\pgfqpoint{2.562770in}{2.184321in}}%
\pgfpathlineto{\pgfqpoint{2.566892in}{2.185843in}}%
\pgfpathlineto{\pgfqpoint{2.571012in}{2.185954in}}%
\pgfpathlineto{\pgfqpoint{2.575135in}{2.184675in}}%
\pgfpathlineto{\pgfqpoint{2.579255in}{2.182024in}}%
\pgfpathlineto{\pgfqpoint{2.583377in}{2.178022in}}%
\pgfpathlineto{\pgfqpoint{2.587497in}{2.172686in}}%
\pgfpathlineto{\pgfqpoint{2.595739in}{2.158089in}}%
\pgfpathlineto{\pgfqpoint{2.603982in}{2.138379in}}%
\pgfpathlineto{\pgfqpoint{2.612224in}{2.113698in}}%
\pgfpathlineto{\pgfqpoint{2.620467in}{2.084186in}}%
\pgfpathlineto{\pgfqpoint{2.628709in}{2.049977in}}%
\pgfpathlineto{\pgfqpoint{2.641074in}{1.990146in}}%
\pgfpathlineto{\pgfqpoint{2.653436in}{1.920477in}}%
\pgfpathlineto{\pgfqpoint{2.665801in}{1.841387in}}%
\pgfpathlineto{\pgfqpoint{2.687071in}{1.697561in}}%
\pgfpathlineto{\pgfqpoint{2.699436in}{1.628725in}}%
\pgfpathlineto{\pgfqpoint{2.707678in}{1.589543in}}%
\pgfpathlineto{\pgfqpoint{2.715921in}{1.555562in}}%
\pgfpathlineto{\pgfqpoint{2.724163in}{1.526641in}}%
\pgfpathlineto{\pgfqpoint{2.732405in}{1.502646in}}%
\pgfpathlineto{\pgfqpoint{2.740648in}{1.483444in}}%
\pgfpathlineto{\pgfqpoint{2.748890in}{1.468906in}}%
\pgfpathlineto{\pgfqpoint{2.757133in}{1.458904in}}%
\pgfpathlineto{\pgfqpoint{2.761253in}{1.455567in}}%
\pgfpathlineto{\pgfqpoint{2.765375in}{1.453318in}}%
\pgfpathlineto{\pgfqpoint{2.769495in}{1.452143in}}%
\pgfpathlineto{\pgfqpoint{2.773618in}{1.452026in}}%
\pgfpathlineto{\pgfqpoint{2.777737in}{1.452954in}}%
\pgfpathlineto{\pgfqpoint{2.781860in}{1.454912in}}%
\pgfpathlineto{\pgfqpoint{2.785980in}{1.457886in}}%
\pgfpathlineto{\pgfqpoint{2.794222in}{1.466825in}}%
\pgfpathlineto{\pgfqpoint{2.802465in}{1.479661in}}%
\pgfpathlineto{\pgfqpoint{2.810707in}{1.496288in}}%
\pgfpathlineto{\pgfqpoint{2.818950in}{1.516601in}}%
\pgfpathlineto{\pgfqpoint{2.827192in}{1.540497in}}%
\pgfpathlineto{\pgfqpoint{2.839557in}{1.582843in}}%
\pgfpathlineto{\pgfqpoint{2.851919in}{1.632703in}}%
\pgfpathlineto{\pgfqpoint{2.864284in}{1.689764in}}%
\pgfpathlineto{\pgfqpoint{2.876647in}{1.753722in}}%
\pgfpathlineto{\pgfqpoint{2.893131in}{1.849228in}}%
\pgfpathlineto{\pgfqpoint{2.910980in}{1.964796in}}%
\pgfpathlineto{\pgfqpoint{2.921317in}{2.026900in}}%
\pgfpathlineto{\pgfqpoint{2.929559in}{2.068699in}}%
\pgfpathlineto{\pgfqpoint{2.937801in}{2.103841in}}%
\pgfpathlineto{\pgfqpoint{2.946044in}{2.132505in}}%
\pgfpathlineto{\pgfqpoint{2.954286in}{2.154869in}}%
\pgfpathlineto{\pgfqpoint{2.962529in}{2.171103in}}%
\pgfpathlineto{\pgfqpoint{2.966649in}{2.176973in}}%
\pgfpathlineto{\pgfqpoint{2.970771in}{2.181374in}}%
\pgfpathlineto{\pgfqpoint{2.974891in}{2.184325in}}%
\pgfpathlineto{\pgfqpoint{2.979013in}{2.185846in}}%
\pgfpathlineto{\pgfqpoint{2.983133in}{2.185957in}}%
\pgfpathlineto{\pgfqpoint{2.987256in}{2.184678in}}%
\pgfpathlineto{\pgfqpoint{2.991376in}{2.182028in}}%
\pgfpathlineto{\pgfqpoint{2.995498in}{2.178025in}}%
\pgfpathlineto{\pgfqpoint{2.999618in}{2.172689in}}%
\pgfpathlineto{\pgfqpoint{3.007861in}{2.158092in}}%
\pgfpathlineto{\pgfqpoint{3.016103in}{2.138382in}}%
\pgfpathlineto{\pgfqpoint{3.024345in}{2.113701in}}%
\pgfpathlineto{\pgfqpoint{3.032588in}{2.084189in}}%
\pgfpathlineto{\pgfqpoint{3.040830in}{2.049980in}}%
\pgfpathlineto{\pgfqpoint{3.053195in}{1.990149in}}%
\pgfpathlineto{\pgfqpoint{3.065558in}{1.920480in}}%
\pgfpathlineto{\pgfqpoint{3.077923in}{1.841390in}}%
\pgfpathlineto{\pgfqpoint{3.099192in}{1.697565in}}%
\pgfpathlineto{\pgfqpoint{3.111557in}{1.628728in}}%
\pgfpathlineto{\pgfqpoint{3.119799in}{1.589547in}}%
\pgfpathlineto{\pgfqpoint{3.128042in}{1.555565in}}%
\pgfpathlineto{\pgfqpoint{3.136284in}{1.526645in}}%
\pgfpathlineto{\pgfqpoint{3.144527in}{1.502650in}}%
\pgfpathlineto{\pgfqpoint{3.152769in}{1.483448in}}%
\pgfpathlineto{\pgfqpoint{3.161011in}{1.468909in}}%
\pgfpathlineto{\pgfqpoint{3.169254in}{1.458908in}}%
\pgfpathlineto{\pgfqpoint{3.173374in}{1.455570in}}%
\pgfpathlineto{\pgfqpoint{3.177496in}{1.453321in}}%
\pgfpathlineto{\pgfqpoint{3.181616in}{1.452146in}}%
\pgfpathlineto{\pgfqpoint{3.185739in}{1.452030in}}%
\pgfpathlineto{\pgfqpoint{3.189859in}{1.452957in}}%
\pgfpathlineto{\pgfqpoint{3.193981in}{1.454915in}}%
\pgfpathlineto{\pgfqpoint{3.198101in}{1.457889in}}%
\pgfpathlineto{\pgfqpoint{3.206343in}{1.466828in}}%
\pgfpathlineto{\pgfqpoint{3.214586in}{1.479665in}}%
\pgfpathlineto{\pgfqpoint{3.222828in}{1.496292in}}%
\pgfpathlineto{\pgfqpoint{3.231071in}{1.516604in}}%
\pgfpathlineto{\pgfqpoint{3.239313in}{1.540500in}}%
\pgfpathlineto{\pgfqpoint{3.251678in}{1.582847in}}%
\pgfpathlineto{\pgfqpoint{3.264040in}{1.632707in}}%
\pgfpathlineto{\pgfqpoint{3.276405in}{1.689767in}}%
\pgfpathlineto{\pgfqpoint{3.288768in}{1.753725in}}%
\pgfpathlineto{\pgfqpoint{3.305253in}{1.849231in}}%
\pgfpathlineto{\pgfqpoint{3.320983in}{1.950708in}}%
\pgfpathlineto{\pgfqpoint{3.320983in}{1.950708in}}%
\pgfusepath{stroke}%
\end{pgfscope}%
\begin{pgfscope}%
\pgfsetrectcap%
\pgfsetmiterjoin%
\pgfsetlinewidth{0.803000pt}%
\definecolor{currentstroke}{rgb}{0.000000,0.000000,0.000000}%
\pgfsetstrokecolor{currentstroke}%
\pgfsetdash{}{0pt}%
\pgfpathmoveto{\pgfqpoint{0.724619in}{1.415315in}}%
\pgfpathlineto{\pgfqpoint{0.724619in}{2.222654in}}%
\pgfusepath{stroke}%
\end{pgfscope}%
\begin{pgfscope}%
\pgfsetrectcap%
\pgfsetmiterjoin%
\pgfsetlinewidth{0.803000pt}%
\definecolor{currentstroke}{rgb}{0.000000,0.000000,0.000000}%
\pgfsetstrokecolor{currentstroke}%
\pgfsetdash{}{0pt}%
\pgfpathmoveto{\pgfqpoint{3.444619in}{1.415315in}}%
\pgfpathlineto{\pgfqpoint{3.444619in}{2.222654in}}%
\pgfusepath{stroke}%
\end{pgfscope}%
\begin{pgfscope}%
\pgfsetrectcap%
\pgfsetmiterjoin%
\pgfsetlinewidth{0.803000pt}%
\definecolor{currentstroke}{rgb}{0.000000,0.000000,0.000000}%
\pgfsetstrokecolor{currentstroke}%
\pgfsetdash{}{0pt}%
\pgfpathmoveto{\pgfqpoint{0.724619in}{1.415315in}}%
\pgfpathlineto{\pgfqpoint{3.444619in}{1.415315in}}%
\pgfusepath{stroke}%
\end{pgfscope}%
\begin{pgfscope}%
\pgfsetrectcap%
\pgfsetmiterjoin%
\pgfsetlinewidth{0.803000pt}%
\definecolor{currentstroke}{rgb}{0.000000,0.000000,0.000000}%
\pgfsetstrokecolor{currentstroke}%
\pgfsetdash{}{0pt}%
\pgfpathmoveto{\pgfqpoint{0.724619in}{2.222654in}}%
\pgfpathlineto{\pgfqpoint{3.444619in}{2.222654in}}%
\pgfusepath{stroke}%
\end{pgfscope}%
\begin{pgfscope}%
\definecolor{textcolor}{rgb}{0.000000,0.000000,0.000000}%
\pgfsetstrokecolor{textcolor}%
\pgfsetfillcolor{textcolor}%
\pgftext[x=3.254219in,y=2.166141in,left,top]{\color{textcolor}{\rmfamily\fontsize{8.000000}{9.600000}\selectfont\catcode`\^=\active\def^{\ifmmode\sp\else\^{}\fi}\catcode`\%=\active\def%{\%}(a)}}%
\end{pgfscope}%
\begin{pgfscope}%
\pgfsetbuttcap%
\pgfsetmiterjoin%
\definecolor{currentfill}{rgb}{1.000000,1.000000,1.000000}%
\pgfsetfillcolor{currentfill}%
\pgfsetlinewidth{0.000000pt}%
\definecolor{currentstroke}{rgb}{0.000000,0.000000,0.000000}%
\pgfsetstrokecolor{currentstroke}%
\pgfsetstrokeopacity{0.000000}%
\pgfsetdash{}{0pt}%
\pgfpathmoveto{\pgfqpoint{0.724619in}{0.462654in}}%
\pgfpathlineto{\pgfqpoint{3.444619in}{0.462654in}}%
\pgfpathlineto{\pgfqpoint{3.444619in}{1.269994in}}%
\pgfpathlineto{\pgfqpoint{0.724619in}{1.269994in}}%
\pgfpathlineto{\pgfqpoint{0.724619in}{0.462654in}}%
\pgfpathclose%
\pgfusepath{fill}%
\end{pgfscope}%
\begin{pgfscope}%
\pgfpathrectangle{\pgfqpoint{0.724619in}{0.462654in}}{\pgfqpoint{2.720000in}{0.807339in}}%
\pgfusepath{clip}%
\pgfsetbuttcap%
\pgfsetroundjoin%
\pgfsetlinewidth{0.501875pt}%
\definecolor{currentstroke}{rgb}{0.690196,0.690196,0.690196}%
\pgfsetstrokecolor{currentstroke}%
\pgfsetstrokeopacity{0.250000}%
\pgfsetdash{{1.850000pt}{0.800000pt}}{0.000000pt}%
\pgfpathmoveto{\pgfqpoint{1.079044in}{0.462654in}}%
\pgfpathlineto{\pgfqpoint{1.079044in}{1.269994in}}%
\pgfusepath{stroke}%
\end{pgfscope}%
\begin{pgfscope}%
\pgfsetbuttcap%
\pgfsetroundjoin%
\definecolor{currentfill}{rgb}{0.000000,0.000000,0.000000}%
\pgfsetfillcolor{currentfill}%
\pgfsetlinewidth{0.803000pt}%
\definecolor{currentstroke}{rgb}{0.000000,0.000000,0.000000}%
\pgfsetstrokecolor{currentstroke}%
\pgfsetdash{}{0pt}%
\pgfsys@defobject{currentmarker}{\pgfqpoint{0.000000in}{-0.048611in}}{\pgfqpoint{0.000000in}{0.000000in}}{%
\pgfpathmoveto{\pgfqpoint{0.000000in}{0.000000in}}%
\pgfpathlineto{\pgfqpoint{0.000000in}{-0.048611in}}%
\pgfusepath{stroke,fill}%
}%
\begin{pgfscope}%
\pgfsys@transformshift{1.079044in}{0.462654in}%
\pgfsys@useobject{currentmarker}{}%
\end{pgfscope}%
\end{pgfscope}%
\begin{pgfscope}%
\definecolor{textcolor}{rgb}{0.000000,0.000000,0.000000}%
\pgfsetstrokecolor{textcolor}%
\pgfsetfillcolor{textcolor}%
\pgftext[x=1.079044in,y=0.365432in,,top]{\color{textcolor}{\rmfamily\fontsize{8.000000}{9.600000}\selectfont\catcode`\^=\active\def^{\ifmmode\sp\else\^{}\fi}\catcode`\%=\active\def%{\%}$\mathdefault{78.04}$}}%
\end{pgfscope}%
\begin{pgfscope}%
\pgfpathrectangle{\pgfqpoint{0.724619in}{0.462654in}}{\pgfqpoint{2.720000in}{0.807339in}}%
\pgfusepath{clip}%
\pgfsetbuttcap%
\pgfsetroundjoin%
\pgfsetlinewidth{0.501875pt}%
\definecolor{currentstroke}{rgb}{0.690196,0.690196,0.690196}%
\pgfsetstrokecolor{currentstroke}%
\pgfsetstrokeopacity{0.250000}%
\pgfsetdash{{1.850000pt}{0.800000pt}}{0.000000pt}%
\pgfpathmoveto{\pgfqpoint{1.606559in}{0.462654in}}%
\pgfpathlineto{\pgfqpoint{1.606559in}{1.269994in}}%
\pgfusepath{stroke}%
\end{pgfscope}%
\begin{pgfscope}%
\pgfsetbuttcap%
\pgfsetroundjoin%
\definecolor{currentfill}{rgb}{0.000000,0.000000,0.000000}%
\pgfsetfillcolor{currentfill}%
\pgfsetlinewidth{0.803000pt}%
\definecolor{currentstroke}{rgb}{0.000000,0.000000,0.000000}%
\pgfsetstrokecolor{currentstroke}%
\pgfsetdash{}{0pt}%
\pgfsys@defobject{currentmarker}{\pgfqpoint{0.000000in}{-0.048611in}}{\pgfqpoint{0.000000in}{0.000000in}}{%
\pgfpathmoveto{\pgfqpoint{0.000000in}{0.000000in}}%
\pgfpathlineto{\pgfqpoint{0.000000in}{-0.048611in}}%
\pgfusepath{stroke,fill}%
}%
\begin{pgfscope}%
\pgfsys@transformshift{1.606559in}{0.462654in}%
\pgfsys@useobject{currentmarker}{}%
\end{pgfscope}%
\end{pgfscope}%
\begin{pgfscope}%
\definecolor{textcolor}{rgb}{0.000000,0.000000,0.000000}%
\pgfsetstrokecolor{textcolor}%
\pgfsetfillcolor{textcolor}%
\pgftext[x=1.606559in,y=0.365432in,,top]{\color{textcolor}{\rmfamily\fontsize{8.000000}{9.600000}\selectfont\catcode`\^=\active\def^{\ifmmode\sp\else\^{}\fi}\catcode`\%=\active\def%{\%}$\mathdefault{78.06}$}}%
\end{pgfscope}%
\begin{pgfscope}%
\pgfpathrectangle{\pgfqpoint{0.724619in}{0.462654in}}{\pgfqpoint{2.720000in}{0.807339in}}%
\pgfusepath{clip}%
\pgfsetbuttcap%
\pgfsetroundjoin%
\pgfsetlinewidth{0.501875pt}%
\definecolor{currentstroke}{rgb}{0.690196,0.690196,0.690196}%
\pgfsetstrokecolor{currentstroke}%
\pgfsetstrokeopacity{0.250000}%
\pgfsetdash{{1.850000pt}{0.800000pt}}{0.000000pt}%
\pgfpathmoveto{\pgfqpoint{2.134074in}{0.462654in}}%
\pgfpathlineto{\pgfqpoint{2.134074in}{1.269994in}}%
\pgfusepath{stroke}%
\end{pgfscope}%
\begin{pgfscope}%
\pgfsetbuttcap%
\pgfsetroundjoin%
\definecolor{currentfill}{rgb}{0.000000,0.000000,0.000000}%
\pgfsetfillcolor{currentfill}%
\pgfsetlinewidth{0.803000pt}%
\definecolor{currentstroke}{rgb}{0.000000,0.000000,0.000000}%
\pgfsetstrokecolor{currentstroke}%
\pgfsetdash{}{0pt}%
\pgfsys@defobject{currentmarker}{\pgfqpoint{0.000000in}{-0.048611in}}{\pgfqpoint{0.000000in}{0.000000in}}{%
\pgfpathmoveto{\pgfqpoint{0.000000in}{0.000000in}}%
\pgfpathlineto{\pgfqpoint{0.000000in}{-0.048611in}}%
\pgfusepath{stroke,fill}%
}%
\begin{pgfscope}%
\pgfsys@transformshift{2.134074in}{0.462654in}%
\pgfsys@useobject{currentmarker}{}%
\end{pgfscope}%
\end{pgfscope}%
\begin{pgfscope}%
\definecolor{textcolor}{rgb}{0.000000,0.000000,0.000000}%
\pgfsetstrokecolor{textcolor}%
\pgfsetfillcolor{textcolor}%
\pgftext[x=2.134074in,y=0.365432in,,top]{\color{textcolor}{\rmfamily\fontsize{8.000000}{9.600000}\selectfont\catcode`\^=\active\def^{\ifmmode\sp\else\^{}\fi}\catcode`\%=\active\def%{\%}$\mathdefault{78.08}$}}%
\end{pgfscope}%
\begin{pgfscope}%
\pgfpathrectangle{\pgfqpoint{0.724619in}{0.462654in}}{\pgfqpoint{2.720000in}{0.807339in}}%
\pgfusepath{clip}%
\pgfsetbuttcap%
\pgfsetroundjoin%
\pgfsetlinewidth{0.501875pt}%
\definecolor{currentstroke}{rgb}{0.690196,0.690196,0.690196}%
\pgfsetstrokecolor{currentstroke}%
\pgfsetstrokeopacity{0.250000}%
\pgfsetdash{{1.850000pt}{0.800000pt}}{0.000000pt}%
\pgfpathmoveto{\pgfqpoint{2.661589in}{0.462654in}}%
\pgfpathlineto{\pgfqpoint{2.661589in}{1.269994in}}%
\pgfusepath{stroke}%
\end{pgfscope}%
\begin{pgfscope}%
\pgfsetbuttcap%
\pgfsetroundjoin%
\definecolor{currentfill}{rgb}{0.000000,0.000000,0.000000}%
\pgfsetfillcolor{currentfill}%
\pgfsetlinewidth{0.803000pt}%
\definecolor{currentstroke}{rgb}{0.000000,0.000000,0.000000}%
\pgfsetstrokecolor{currentstroke}%
\pgfsetdash{}{0pt}%
\pgfsys@defobject{currentmarker}{\pgfqpoint{0.000000in}{-0.048611in}}{\pgfqpoint{0.000000in}{0.000000in}}{%
\pgfpathmoveto{\pgfqpoint{0.000000in}{0.000000in}}%
\pgfpathlineto{\pgfqpoint{0.000000in}{-0.048611in}}%
\pgfusepath{stroke,fill}%
}%
\begin{pgfscope}%
\pgfsys@transformshift{2.661589in}{0.462654in}%
\pgfsys@useobject{currentmarker}{}%
\end{pgfscope}%
\end{pgfscope}%
\begin{pgfscope}%
\definecolor{textcolor}{rgb}{0.000000,0.000000,0.000000}%
\pgfsetstrokecolor{textcolor}%
\pgfsetfillcolor{textcolor}%
\pgftext[x=2.661589in,y=0.365432in,,top]{\color{textcolor}{\rmfamily\fontsize{8.000000}{9.600000}\selectfont\catcode`\^=\active\def^{\ifmmode\sp\else\^{}\fi}\catcode`\%=\active\def%{\%}$\mathdefault{78.10}$}}%
\end{pgfscope}%
\begin{pgfscope}%
\pgfpathrectangle{\pgfqpoint{0.724619in}{0.462654in}}{\pgfqpoint{2.720000in}{0.807339in}}%
\pgfusepath{clip}%
\pgfsetbuttcap%
\pgfsetroundjoin%
\pgfsetlinewidth{0.501875pt}%
\definecolor{currentstroke}{rgb}{0.690196,0.690196,0.690196}%
\pgfsetstrokecolor{currentstroke}%
\pgfsetstrokeopacity{0.250000}%
\pgfsetdash{{1.850000pt}{0.800000pt}}{0.000000pt}%
\pgfpathmoveto{\pgfqpoint{3.189104in}{0.462654in}}%
\pgfpathlineto{\pgfqpoint{3.189104in}{1.269994in}}%
\pgfusepath{stroke}%
\end{pgfscope}%
\begin{pgfscope}%
\pgfsetbuttcap%
\pgfsetroundjoin%
\definecolor{currentfill}{rgb}{0.000000,0.000000,0.000000}%
\pgfsetfillcolor{currentfill}%
\pgfsetlinewidth{0.803000pt}%
\definecolor{currentstroke}{rgb}{0.000000,0.000000,0.000000}%
\pgfsetstrokecolor{currentstroke}%
\pgfsetdash{}{0pt}%
\pgfsys@defobject{currentmarker}{\pgfqpoint{0.000000in}{-0.048611in}}{\pgfqpoint{0.000000in}{0.000000in}}{%
\pgfpathmoveto{\pgfqpoint{0.000000in}{0.000000in}}%
\pgfpathlineto{\pgfqpoint{0.000000in}{-0.048611in}}%
\pgfusepath{stroke,fill}%
}%
\begin{pgfscope}%
\pgfsys@transformshift{3.189104in}{0.462654in}%
\pgfsys@useobject{currentmarker}{}%
\end{pgfscope}%
\end{pgfscope}%
\begin{pgfscope}%
\definecolor{textcolor}{rgb}{0.000000,0.000000,0.000000}%
\pgfsetstrokecolor{textcolor}%
\pgfsetfillcolor{textcolor}%
\pgftext[x=3.189104in,y=0.365432in,,top]{\color{textcolor}{\rmfamily\fontsize{8.000000}{9.600000}\selectfont\catcode`\^=\active\def^{\ifmmode\sp\else\^{}\fi}\catcode`\%=\active\def%{\%}$\mathdefault{78.12}$}}%
\end{pgfscope}%
\begin{pgfscope}%
\definecolor{textcolor}{rgb}{0.000000,0.000000,0.000000}%
\pgfsetstrokecolor{textcolor}%
\pgfsetfillcolor{textcolor}%
\pgftext[x=2.084619in,y=0.211111in,,top]{\color{textcolor}{\rmfamily\fontsize{8.000000}{9.600000}\selectfont\catcode`\^=\active\def^{\ifmmode\sp\else\^{}\fi}\catcode`\%=\active\def%{\%}Time (ms)}}%
\end{pgfscope}%
\begin{pgfscope}%
\pgfpathrectangle{\pgfqpoint{0.724619in}{0.462654in}}{\pgfqpoint{2.720000in}{0.807339in}}%
\pgfusepath{clip}%
\pgfsetbuttcap%
\pgfsetroundjoin%
\pgfsetlinewidth{0.501875pt}%
\definecolor{currentstroke}{rgb}{0.690196,0.690196,0.690196}%
\pgfsetstrokecolor{currentstroke}%
\pgfsetstrokeopacity{0.250000}%
\pgfsetdash{{1.850000pt}{0.800000pt}}{0.000000pt}%
\pgfpathmoveto{\pgfqpoint{0.724619in}{0.548442in}}%
\pgfpathlineto{\pgfqpoint{3.444619in}{0.548442in}}%
\pgfusepath{stroke}%
\end{pgfscope}%
\begin{pgfscope}%
\pgfsetbuttcap%
\pgfsetroundjoin%
\definecolor{currentfill}{rgb}{0.000000,0.000000,0.000000}%
\pgfsetfillcolor{currentfill}%
\pgfsetlinewidth{0.803000pt}%
\definecolor{currentstroke}{rgb}{0.000000,0.000000,0.000000}%
\pgfsetstrokecolor{currentstroke}%
\pgfsetdash{}{0pt}%
\pgfsys@defobject{currentmarker}{\pgfqpoint{-0.048611in}{0.000000in}}{\pgfqpoint{-0.000000in}{0.000000in}}{%
\pgfpathmoveto{\pgfqpoint{-0.000000in}{0.000000in}}%
\pgfpathlineto{\pgfqpoint{-0.048611in}{0.000000in}}%
\pgfusepath{stroke,fill}%
}%
\begin{pgfscope}%
\pgfsys@transformshift{0.724619in}{0.548442in}%
\pgfsys@useobject{currentmarker}{}%
\end{pgfscope}%
\end{pgfscope}%
\begin{pgfscope}%
\definecolor{textcolor}{rgb}{0.000000,0.000000,0.000000}%
\pgfsetstrokecolor{textcolor}%
\pgfsetfillcolor{textcolor}%
\pgftext[x=0.417518in, y=0.509861in, left, base]{\color{textcolor}{\rmfamily\fontsize{8.000000}{9.600000}\selectfont\catcode`\^=\active\def^{\ifmmode\sp\else\^{}\fi}\catcode`\%=\active\def%{\%}$\mathdefault{2.97}$}}%
\end{pgfscope}%
\begin{pgfscope}%
\pgfpathrectangle{\pgfqpoint{0.724619in}{0.462654in}}{\pgfqpoint{2.720000in}{0.807339in}}%
\pgfusepath{clip}%
\pgfsetbuttcap%
\pgfsetroundjoin%
\pgfsetlinewidth{0.501875pt}%
\definecolor{currentstroke}{rgb}{0.690196,0.690196,0.690196}%
\pgfsetstrokecolor{currentstroke}%
\pgfsetstrokeopacity{0.250000}%
\pgfsetdash{{1.850000pt}{0.800000pt}}{0.000000pt}%
\pgfpathmoveto{\pgfqpoint{0.724619in}{0.776735in}}%
\pgfpathlineto{\pgfqpoint{3.444619in}{0.776735in}}%
\pgfusepath{stroke}%
\end{pgfscope}%
\begin{pgfscope}%
\pgfsetbuttcap%
\pgfsetroundjoin%
\definecolor{currentfill}{rgb}{0.000000,0.000000,0.000000}%
\pgfsetfillcolor{currentfill}%
\pgfsetlinewidth{0.803000pt}%
\definecolor{currentstroke}{rgb}{0.000000,0.000000,0.000000}%
\pgfsetstrokecolor{currentstroke}%
\pgfsetdash{}{0pt}%
\pgfsys@defobject{currentmarker}{\pgfqpoint{-0.048611in}{0.000000in}}{\pgfqpoint{-0.000000in}{0.000000in}}{%
\pgfpathmoveto{\pgfqpoint{-0.000000in}{0.000000in}}%
\pgfpathlineto{\pgfqpoint{-0.048611in}{0.000000in}}%
\pgfusepath{stroke,fill}%
}%
\begin{pgfscope}%
\pgfsys@transformshift{0.724619in}{0.776735in}%
\pgfsys@useobject{currentmarker}{}%
\end{pgfscope}%
\end{pgfscope}%
\begin{pgfscope}%
\definecolor{textcolor}{rgb}{0.000000,0.000000,0.000000}%
\pgfsetstrokecolor{textcolor}%
\pgfsetfillcolor{textcolor}%
\pgftext[x=0.417518in, y=0.738155in, left, base]{\color{textcolor}{\rmfamily\fontsize{8.000000}{9.600000}\selectfont\catcode`\^=\active\def^{\ifmmode\sp\else\^{}\fi}\catcode`\%=\active\def%{\%}$\mathdefault{2.98}$}}%
\end{pgfscope}%
\begin{pgfscope}%
\pgfpathrectangle{\pgfqpoint{0.724619in}{0.462654in}}{\pgfqpoint{2.720000in}{0.807339in}}%
\pgfusepath{clip}%
\pgfsetbuttcap%
\pgfsetroundjoin%
\pgfsetlinewidth{0.501875pt}%
\definecolor{currentstroke}{rgb}{0.690196,0.690196,0.690196}%
\pgfsetstrokecolor{currentstroke}%
\pgfsetstrokeopacity{0.250000}%
\pgfsetdash{{1.850000pt}{0.800000pt}}{0.000000pt}%
\pgfpathmoveto{\pgfqpoint{0.724619in}{1.005029in}}%
\pgfpathlineto{\pgfqpoint{3.444619in}{1.005029in}}%
\pgfusepath{stroke}%
\end{pgfscope}%
\begin{pgfscope}%
\pgfsetbuttcap%
\pgfsetroundjoin%
\definecolor{currentfill}{rgb}{0.000000,0.000000,0.000000}%
\pgfsetfillcolor{currentfill}%
\pgfsetlinewidth{0.803000pt}%
\definecolor{currentstroke}{rgb}{0.000000,0.000000,0.000000}%
\pgfsetstrokecolor{currentstroke}%
\pgfsetdash{}{0pt}%
\pgfsys@defobject{currentmarker}{\pgfqpoint{-0.048611in}{0.000000in}}{\pgfqpoint{-0.000000in}{0.000000in}}{%
\pgfpathmoveto{\pgfqpoint{-0.000000in}{0.000000in}}%
\pgfpathlineto{\pgfqpoint{-0.048611in}{0.000000in}}%
\pgfusepath{stroke,fill}%
}%
\begin{pgfscope}%
\pgfsys@transformshift{0.724619in}{1.005029in}%
\pgfsys@useobject{currentmarker}{}%
\end{pgfscope}%
\end{pgfscope}%
\begin{pgfscope}%
\definecolor{textcolor}{rgb}{0.000000,0.000000,0.000000}%
\pgfsetstrokecolor{textcolor}%
\pgfsetfillcolor{textcolor}%
\pgftext[x=0.417518in, y=0.966449in, left, base]{\color{textcolor}{\rmfamily\fontsize{8.000000}{9.600000}\selectfont\catcode`\^=\active\def^{\ifmmode\sp\else\^{}\fi}\catcode`\%=\active\def%{\%}$\mathdefault{2.99}$}}%
\end{pgfscope}%
\begin{pgfscope}%
\pgfpathrectangle{\pgfqpoint{0.724619in}{0.462654in}}{\pgfqpoint{2.720000in}{0.807339in}}%
\pgfusepath{clip}%
\pgfsetbuttcap%
\pgfsetroundjoin%
\pgfsetlinewidth{0.501875pt}%
\definecolor{currentstroke}{rgb}{0.690196,0.690196,0.690196}%
\pgfsetstrokecolor{currentstroke}%
\pgfsetstrokeopacity{0.250000}%
\pgfsetdash{{1.850000pt}{0.800000pt}}{0.000000pt}%
\pgfpathmoveto{\pgfqpoint{0.724619in}{1.233323in}}%
\pgfpathlineto{\pgfqpoint{3.444619in}{1.233323in}}%
\pgfusepath{stroke}%
\end{pgfscope}%
\begin{pgfscope}%
\pgfsetbuttcap%
\pgfsetroundjoin%
\definecolor{currentfill}{rgb}{0.000000,0.000000,0.000000}%
\pgfsetfillcolor{currentfill}%
\pgfsetlinewidth{0.803000pt}%
\definecolor{currentstroke}{rgb}{0.000000,0.000000,0.000000}%
\pgfsetstrokecolor{currentstroke}%
\pgfsetdash{}{0pt}%
\pgfsys@defobject{currentmarker}{\pgfqpoint{-0.048611in}{0.000000in}}{\pgfqpoint{-0.000000in}{0.000000in}}{%
\pgfpathmoveto{\pgfqpoint{-0.000000in}{0.000000in}}%
\pgfpathlineto{\pgfqpoint{-0.048611in}{0.000000in}}%
\pgfusepath{stroke,fill}%
}%
\begin{pgfscope}%
\pgfsys@transformshift{0.724619in}{1.233323in}%
\pgfsys@useobject{currentmarker}{}%
\end{pgfscope}%
\end{pgfscope}%
\begin{pgfscope}%
\definecolor{textcolor}{rgb}{0.000000,0.000000,0.000000}%
\pgfsetstrokecolor{textcolor}%
\pgfsetfillcolor{textcolor}%
\pgftext[x=0.417518in, y=1.194743in, left, base]{\color{textcolor}{\rmfamily\fontsize{8.000000}{9.600000}\selectfont\catcode`\^=\active\def^{\ifmmode\sp\else\^{}\fi}\catcode`\%=\active\def%{\%}$\mathdefault{3.00}$}}%
\end{pgfscope}%
\begin{pgfscope}%
\definecolor{textcolor}{rgb}{0.000000,0.000000,0.000000}%
\pgfsetstrokecolor{textcolor}%
\pgfsetfillcolor{textcolor}%
\pgftext[x=0.361962in,y=0.866324in,,bottom,rotate=90.000000]{\color{textcolor}{\rmfamily\fontsize{8.000000}{9.600000}\selectfont\catcode`\^=\active\def^{\ifmmode\sp\else\^{}\fi}\catcode`\%=\active\def%{\%}$i_L$ (A)}}%
\end{pgfscope}%
\begin{pgfscope}%
\pgfpathrectangle{\pgfqpoint{0.724619in}{0.462654in}}{\pgfqpoint{2.720000in}{0.807339in}}%
\pgfusepath{clip}%
\pgfsetrectcap%
\pgfsetroundjoin%
\pgfsetlinewidth{1.505625pt}%
\definecolor{currentstroke}{rgb}{0.835294,0.368627,0.000000}%
\pgfsetstrokecolor{currentstroke}%
\pgfsetdash{}{0pt}%
\pgfpathmoveto{\pgfqpoint{0.848256in}{0.499395in}}%
\pgfpathlineto{\pgfqpoint{0.848266in}{0.499358in}}%
\pgfpathlineto{\pgfqpoint{0.848543in}{0.500494in}}%
\pgfpathlineto{\pgfqpoint{1.024931in}{1.233297in}}%
\pgfpathlineto{\pgfqpoint{1.025111in}{1.232736in}}%
\pgfpathlineto{\pgfqpoint{1.244647in}{0.548592in}}%
\pgfpathlineto{\pgfqpoint{1.260388in}{0.499357in}}%
\pgfpathlineto{\pgfqpoint{1.260665in}{0.500492in}}%
\pgfpathlineto{\pgfqpoint{1.437052in}{1.233295in}}%
\pgfpathlineto{\pgfqpoint{1.437232in}{1.232734in}}%
\pgfpathlineto{\pgfqpoint{1.656768in}{0.548591in}}%
\pgfpathlineto{\pgfqpoint{1.672509in}{0.499355in}}%
\pgfpathlineto{\pgfqpoint{1.672786in}{0.500490in}}%
\pgfpathlineto{\pgfqpoint{1.849174in}{1.233294in}}%
\pgfpathlineto{\pgfqpoint{1.849353in}{1.232733in}}%
\pgfpathlineto{\pgfqpoint{2.084630in}{0.499354in}}%
\pgfpathlineto{\pgfqpoint{2.085168in}{0.501576in}}%
\pgfpathlineto{\pgfqpoint{2.261295in}{1.233292in}}%
\pgfpathlineto{\pgfqpoint{2.261474in}{1.232731in}}%
\pgfpathlineto{\pgfqpoint{2.496751in}{0.499353in}}%
\pgfpathlineto{\pgfqpoint{2.497289in}{0.501575in}}%
\pgfpathlineto{\pgfqpoint{2.673416in}{1.233291in}}%
\pgfpathlineto{\pgfqpoint{2.673595in}{1.232730in}}%
\pgfpathlineto{\pgfqpoint{2.893131in}{0.548587in}}%
\pgfpathlineto{\pgfqpoint{2.908872in}{0.499352in}}%
\pgfpathlineto{\pgfqpoint{2.909149in}{0.500487in}}%
\pgfpathlineto{\pgfqpoint{3.085537in}{1.233290in}}%
\pgfpathlineto{\pgfqpoint{3.085717in}{1.232729in}}%
\pgfpathlineto{\pgfqpoint{3.305253in}{0.548586in}}%
\pgfpathlineto{\pgfqpoint{3.320983in}{0.499387in}}%
\pgfpathlineto{\pgfqpoint{3.320983in}{0.499387in}}%
\pgfusepath{stroke}%
\end{pgfscope}%
\begin{pgfscope}%
\pgfsetrectcap%
\pgfsetmiterjoin%
\pgfsetlinewidth{0.803000pt}%
\definecolor{currentstroke}{rgb}{0.000000,0.000000,0.000000}%
\pgfsetstrokecolor{currentstroke}%
\pgfsetdash{}{0pt}%
\pgfpathmoveto{\pgfqpoint{0.724619in}{0.462654in}}%
\pgfpathlineto{\pgfqpoint{0.724619in}{1.269994in}}%
\pgfusepath{stroke}%
\end{pgfscope}%
\begin{pgfscope}%
\pgfsetrectcap%
\pgfsetmiterjoin%
\pgfsetlinewidth{0.803000pt}%
\definecolor{currentstroke}{rgb}{0.000000,0.000000,0.000000}%
\pgfsetstrokecolor{currentstroke}%
\pgfsetdash{}{0pt}%
\pgfpathmoveto{\pgfqpoint{3.444619in}{0.462654in}}%
\pgfpathlineto{\pgfqpoint{3.444619in}{1.269994in}}%
\pgfusepath{stroke}%
\end{pgfscope}%
\begin{pgfscope}%
\pgfsetrectcap%
\pgfsetmiterjoin%
\pgfsetlinewidth{0.803000pt}%
\definecolor{currentstroke}{rgb}{0.000000,0.000000,0.000000}%
\pgfsetstrokecolor{currentstroke}%
\pgfsetdash{}{0pt}%
\pgfpathmoveto{\pgfqpoint{0.724619in}{0.462654in}}%
\pgfpathlineto{\pgfqpoint{3.444619in}{0.462654in}}%
\pgfusepath{stroke}%
\end{pgfscope}%
\begin{pgfscope}%
\pgfsetrectcap%
\pgfsetmiterjoin%
\pgfsetlinewidth{0.803000pt}%
\definecolor{currentstroke}{rgb}{0.000000,0.000000,0.000000}%
\pgfsetstrokecolor{currentstroke}%
\pgfsetdash{}{0pt}%
\pgfpathmoveto{\pgfqpoint{0.724619in}{1.269994in}}%
\pgfpathlineto{\pgfqpoint{3.444619in}{1.269994in}}%
\pgfusepath{stroke}%
\end{pgfscope}%
\begin{pgfscope}%
\definecolor{textcolor}{rgb}{0.000000,0.000000,0.000000}%
\pgfsetstrokecolor{textcolor}%
\pgfsetfillcolor{textcolor}%
\pgftext[x=0.779019in,y=1.213480in,left,top]{\color{textcolor}{\rmfamily\fontsize{8.000000}{9.600000}\selectfont\catcode`\^=\active\def^{\ifmmode\sp\else\^{}\fi}\catcode`\%=\active\def%{\%}(b)}}%
\end{pgfscope}%
\end{pgfpicture}%
\makeatother%
\endgroup%

	\caption{(a) tensão de saída e (b) corrente do indutor $L_2$ do circuito Cuk.}
	\label{fig:ex3_2}
\end{figure}

\medskip
\textbf{d)} Determine a corrente média de entrada e a ondulação da corrente de saída (pico-a-pico).

\medskip
\textbf{Resolução}: a corrente média de entrada pode ser determinada por
\begin{equation*}
	I_{L1} = \frac{E}{R_o}\big(\frac{D}{1 - D}\big)^2 = 3~\mathrm{A}.
\end{equation*}
Enquanto que o \textit{ripple} da corrente de saída pode ser calculado por
\begin{equation*}
	\Delta i_{o,pico} = \frac{\Delta v_{pico}}{R_o} = 55{,}6~\mathrm{mA}.
\end{equation*}

Estes valores podem ser verificados por meio de simulação como mostrado na Figura~\ref{fig:ex3_3}.
\begin{figure}[htb!]
	\centering
	%% Creator: Matplotlib, PGF backend
%%
%% To include the figure in your LaTeX document, write
%%   \input{<filename>.pgf}
%%
%% Make sure the required packages are loaded in your preamble
%%   \usepackage{pgf}
%%
%% Also ensure that all the required font packages are loaded; for instance,
%% the lmodern package is sometimes necessary when using math font.
%%   \usepackage{lmodern}
%%
%% Figures using additional raster images can only be included by \input if
%% they are in the same directory as the main LaTeX file. For loading figures
%% from other directories you can use the `import` package
%%   \usepackage{import}
%%
%% and then include the figures with
%%   \import{<path to file>}{<filename>.pgf}
%%
%% Matplotlib used the following preamble
%%   \def\mathdefault#1{#1}
%%   \everymath=\expandafter{\the\everymath\displaystyle}
%%   \IfFileExists{scrextend.sty}{
%%     \usepackage[fontsize=8.000000pt]{scrextend}
%%   }{
%%     \renewcommand{\normalsize}{\fontsize{8.000000}{9.600000}\selectfont}
%%     \normalsize
%%   }
%%   
%%           \usepackage{amsmath}
%%           \usepackage{amssymb}
%%           \usepackage{mathtools}
%%           \usepackage{dsfont}
%%           \providecommand{\mathdefault}[1]{#1}
%%       
%%   \makeatletter\@ifpackageloaded{underscore}{}{\usepackage[strings]{underscore}}\makeatother
%%
\begingroup%
\makeatletter%
\begin{pgfpicture}%
\pgfpathrectangle{\pgfpointorigin}{\pgfqpoint{3.544619in}{2.324564in}}%
\pgfusepath{use as bounding box, clip}%
\begin{pgfscope}%
\pgfsetbuttcap%
\pgfsetmiterjoin%
\definecolor{currentfill}{rgb}{1.000000,1.000000,1.000000}%
\pgfsetfillcolor{currentfill}%
\pgfsetlinewidth{0.000000pt}%
\definecolor{currentstroke}{rgb}{1.000000,1.000000,1.000000}%
\pgfsetstrokecolor{currentstroke}%
\pgfsetdash{}{0pt}%
\pgfpathmoveto{\pgfqpoint{0.000000in}{0.000000in}}%
\pgfpathlineto{\pgfqpoint{3.544619in}{0.000000in}}%
\pgfpathlineto{\pgfqpoint{3.544619in}{2.324564in}}%
\pgfpathlineto{\pgfqpoint{0.000000in}{2.324564in}}%
\pgfpathlineto{\pgfqpoint{0.000000in}{0.000000in}}%
\pgfpathclose%
\pgfusepath{fill}%
\end{pgfscope}%
\begin{pgfscope}%
\pgfsetbuttcap%
\pgfsetmiterjoin%
\definecolor{currentfill}{rgb}{1.000000,1.000000,1.000000}%
\pgfsetfillcolor{currentfill}%
\pgfsetlinewidth{0.000000pt}%
\definecolor{currentstroke}{rgb}{0.000000,0.000000,0.000000}%
\pgfsetstrokecolor{currentstroke}%
\pgfsetstrokeopacity{0.000000}%
\pgfsetdash{}{0pt}%
\pgfpathmoveto{\pgfqpoint{0.724619in}{1.415315in}}%
\pgfpathlineto{\pgfqpoint{3.444619in}{1.415315in}}%
\pgfpathlineto{\pgfqpoint{3.444619in}{2.222654in}}%
\pgfpathlineto{\pgfqpoint{0.724619in}{2.222654in}}%
\pgfpathlineto{\pgfqpoint{0.724619in}{1.415315in}}%
\pgfpathclose%
\pgfusepath{fill}%
\end{pgfscope}%
\begin{pgfscope}%
\pgfpathrectangle{\pgfqpoint{0.724619in}{1.415315in}}{\pgfqpoint{2.720000in}{0.807339in}}%
\pgfusepath{clip}%
\pgfsetbuttcap%
\pgfsetroundjoin%
\pgfsetlinewidth{0.501875pt}%
\definecolor{currentstroke}{rgb}{0.690196,0.690196,0.690196}%
\pgfsetstrokecolor{currentstroke}%
\pgfsetstrokeopacity{0.250000}%
\pgfsetdash{{1.850000pt}{0.800000pt}}{0.000000pt}%
\pgfpathmoveto{\pgfqpoint{1.079044in}{1.415315in}}%
\pgfpathlineto{\pgfqpoint{1.079044in}{2.222654in}}%
\pgfusepath{stroke}%
\end{pgfscope}%
\begin{pgfscope}%
\pgfsetbuttcap%
\pgfsetroundjoin%
\definecolor{currentfill}{rgb}{0.000000,0.000000,0.000000}%
\pgfsetfillcolor{currentfill}%
\pgfsetlinewidth{0.803000pt}%
\definecolor{currentstroke}{rgb}{0.000000,0.000000,0.000000}%
\pgfsetstrokecolor{currentstroke}%
\pgfsetdash{}{0pt}%
\pgfsys@defobject{currentmarker}{\pgfqpoint{0.000000in}{-0.048611in}}{\pgfqpoint{0.000000in}{0.000000in}}{%
\pgfpathmoveto{\pgfqpoint{0.000000in}{0.000000in}}%
\pgfpathlineto{\pgfqpoint{0.000000in}{-0.048611in}}%
\pgfusepath{stroke,fill}%
}%
\begin{pgfscope}%
\pgfsys@transformshift{1.079044in}{1.415315in}%
\pgfsys@useobject{currentmarker}{}%
\end{pgfscope}%
\end{pgfscope}%
\begin{pgfscope}%
\pgfpathrectangle{\pgfqpoint{0.724619in}{1.415315in}}{\pgfqpoint{2.720000in}{0.807339in}}%
\pgfusepath{clip}%
\pgfsetbuttcap%
\pgfsetroundjoin%
\pgfsetlinewidth{0.501875pt}%
\definecolor{currentstroke}{rgb}{0.690196,0.690196,0.690196}%
\pgfsetstrokecolor{currentstroke}%
\pgfsetstrokeopacity{0.250000}%
\pgfsetdash{{1.850000pt}{0.800000pt}}{0.000000pt}%
\pgfpathmoveto{\pgfqpoint{1.606559in}{1.415315in}}%
\pgfpathlineto{\pgfqpoint{1.606559in}{2.222654in}}%
\pgfusepath{stroke}%
\end{pgfscope}%
\begin{pgfscope}%
\pgfsetbuttcap%
\pgfsetroundjoin%
\definecolor{currentfill}{rgb}{0.000000,0.000000,0.000000}%
\pgfsetfillcolor{currentfill}%
\pgfsetlinewidth{0.803000pt}%
\definecolor{currentstroke}{rgb}{0.000000,0.000000,0.000000}%
\pgfsetstrokecolor{currentstroke}%
\pgfsetdash{}{0pt}%
\pgfsys@defobject{currentmarker}{\pgfqpoint{0.000000in}{-0.048611in}}{\pgfqpoint{0.000000in}{0.000000in}}{%
\pgfpathmoveto{\pgfqpoint{0.000000in}{0.000000in}}%
\pgfpathlineto{\pgfqpoint{0.000000in}{-0.048611in}}%
\pgfusepath{stroke,fill}%
}%
\begin{pgfscope}%
\pgfsys@transformshift{1.606559in}{1.415315in}%
\pgfsys@useobject{currentmarker}{}%
\end{pgfscope}%
\end{pgfscope}%
\begin{pgfscope}%
\pgfpathrectangle{\pgfqpoint{0.724619in}{1.415315in}}{\pgfqpoint{2.720000in}{0.807339in}}%
\pgfusepath{clip}%
\pgfsetbuttcap%
\pgfsetroundjoin%
\pgfsetlinewidth{0.501875pt}%
\definecolor{currentstroke}{rgb}{0.690196,0.690196,0.690196}%
\pgfsetstrokecolor{currentstroke}%
\pgfsetstrokeopacity{0.250000}%
\pgfsetdash{{1.850000pt}{0.800000pt}}{0.000000pt}%
\pgfpathmoveto{\pgfqpoint{2.134074in}{1.415315in}}%
\pgfpathlineto{\pgfqpoint{2.134074in}{2.222654in}}%
\pgfusepath{stroke}%
\end{pgfscope}%
\begin{pgfscope}%
\pgfsetbuttcap%
\pgfsetroundjoin%
\definecolor{currentfill}{rgb}{0.000000,0.000000,0.000000}%
\pgfsetfillcolor{currentfill}%
\pgfsetlinewidth{0.803000pt}%
\definecolor{currentstroke}{rgb}{0.000000,0.000000,0.000000}%
\pgfsetstrokecolor{currentstroke}%
\pgfsetdash{}{0pt}%
\pgfsys@defobject{currentmarker}{\pgfqpoint{0.000000in}{-0.048611in}}{\pgfqpoint{0.000000in}{0.000000in}}{%
\pgfpathmoveto{\pgfqpoint{0.000000in}{0.000000in}}%
\pgfpathlineto{\pgfqpoint{0.000000in}{-0.048611in}}%
\pgfusepath{stroke,fill}%
}%
\begin{pgfscope}%
\pgfsys@transformshift{2.134074in}{1.415315in}%
\pgfsys@useobject{currentmarker}{}%
\end{pgfscope}%
\end{pgfscope}%
\begin{pgfscope}%
\pgfpathrectangle{\pgfqpoint{0.724619in}{1.415315in}}{\pgfqpoint{2.720000in}{0.807339in}}%
\pgfusepath{clip}%
\pgfsetbuttcap%
\pgfsetroundjoin%
\pgfsetlinewidth{0.501875pt}%
\definecolor{currentstroke}{rgb}{0.690196,0.690196,0.690196}%
\pgfsetstrokecolor{currentstroke}%
\pgfsetstrokeopacity{0.250000}%
\pgfsetdash{{1.850000pt}{0.800000pt}}{0.000000pt}%
\pgfpathmoveto{\pgfqpoint{2.661589in}{1.415315in}}%
\pgfpathlineto{\pgfqpoint{2.661589in}{2.222654in}}%
\pgfusepath{stroke}%
\end{pgfscope}%
\begin{pgfscope}%
\pgfsetbuttcap%
\pgfsetroundjoin%
\definecolor{currentfill}{rgb}{0.000000,0.000000,0.000000}%
\pgfsetfillcolor{currentfill}%
\pgfsetlinewidth{0.803000pt}%
\definecolor{currentstroke}{rgb}{0.000000,0.000000,0.000000}%
\pgfsetstrokecolor{currentstroke}%
\pgfsetdash{}{0pt}%
\pgfsys@defobject{currentmarker}{\pgfqpoint{0.000000in}{-0.048611in}}{\pgfqpoint{0.000000in}{0.000000in}}{%
\pgfpathmoveto{\pgfqpoint{0.000000in}{0.000000in}}%
\pgfpathlineto{\pgfqpoint{0.000000in}{-0.048611in}}%
\pgfusepath{stroke,fill}%
}%
\begin{pgfscope}%
\pgfsys@transformshift{2.661589in}{1.415315in}%
\pgfsys@useobject{currentmarker}{}%
\end{pgfscope}%
\end{pgfscope}%
\begin{pgfscope}%
\pgfpathrectangle{\pgfqpoint{0.724619in}{1.415315in}}{\pgfqpoint{2.720000in}{0.807339in}}%
\pgfusepath{clip}%
\pgfsetbuttcap%
\pgfsetroundjoin%
\pgfsetlinewidth{0.501875pt}%
\definecolor{currentstroke}{rgb}{0.690196,0.690196,0.690196}%
\pgfsetstrokecolor{currentstroke}%
\pgfsetstrokeopacity{0.250000}%
\pgfsetdash{{1.850000pt}{0.800000pt}}{0.000000pt}%
\pgfpathmoveto{\pgfqpoint{3.189104in}{1.415315in}}%
\pgfpathlineto{\pgfqpoint{3.189104in}{2.222654in}}%
\pgfusepath{stroke}%
\end{pgfscope}%
\begin{pgfscope}%
\pgfsetbuttcap%
\pgfsetroundjoin%
\definecolor{currentfill}{rgb}{0.000000,0.000000,0.000000}%
\pgfsetfillcolor{currentfill}%
\pgfsetlinewidth{0.803000pt}%
\definecolor{currentstroke}{rgb}{0.000000,0.000000,0.000000}%
\pgfsetstrokecolor{currentstroke}%
\pgfsetdash{}{0pt}%
\pgfsys@defobject{currentmarker}{\pgfqpoint{0.000000in}{-0.048611in}}{\pgfqpoint{0.000000in}{0.000000in}}{%
\pgfpathmoveto{\pgfqpoint{0.000000in}{0.000000in}}%
\pgfpathlineto{\pgfqpoint{0.000000in}{-0.048611in}}%
\pgfusepath{stroke,fill}%
}%
\begin{pgfscope}%
\pgfsys@transformshift{3.189104in}{1.415315in}%
\pgfsys@useobject{currentmarker}{}%
\end{pgfscope}%
\end{pgfscope}%
\begin{pgfscope}%
\pgfpathrectangle{\pgfqpoint{0.724619in}{1.415315in}}{\pgfqpoint{2.720000in}{0.807339in}}%
\pgfusepath{clip}%
\pgfsetbuttcap%
\pgfsetroundjoin%
\pgfsetlinewidth{0.501875pt}%
\definecolor{currentstroke}{rgb}{0.690196,0.690196,0.690196}%
\pgfsetstrokecolor{currentstroke}%
\pgfsetstrokeopacity{0.250000}%
\pgfsetdash{{1.850000pt}{0.800000pt}}{0.000000pt}%
\pgfpathmoveto{\pgfqpoint{0.724619in}{1.501102in}}%
\pgfpathlineto{\pgfqpoint{3.444619in}{1.501102in}}%
\pgfusepath{stroke}%
\end{pgfscope}%
\begin{pgfscope}%
\pgfsetbuttcap%
\pgfsetroundjoin%
\definecolor{currentfill}{rgb}{0.000000,0.000000,0.000000}%
\pgfsetfillcolor{currentfill}%
\pgfsetlinewidth{0.803000pt}%
\definecolor{currentstroke}{rgb}{0.000000,0.000000,0.000000}%
\pgfsetstrokecolor{currentstroke}%
\pgfsetdash{}{0pt}%
\pgfsys@defobject{currentmarker}{\pgfqpoint{-0.048611in}{0.000000in}}{\pgfqpoint{-0.000000in}{0.000000in}}{%
\pgfpathmoveto{\pgfqpoint{-0.000000in}{0.000000in}}%
\pgfpathlineto{\pgfqpoint{-0.048611in}{0.000000in}}%
\pgfusepath{stroke,fill}%
}%
\begin{pgfscope}%
\pgfsys@transformshift{0.724619in}{1.501102in}%
\pgfsys@useobject{currentmarker}{}%
\end{pgfscope}%
\end{pgfscope}%
\begin{pgfscope}%
\definecolor{textcolor}{rgb}{0.000000,0.000000,0.000000}%
\pgfsetstrokecolor{textcolor}%
\pgfsetfillcolor{textcolor}%
\pgftext[x=0.417518in, y=1.462522in, left, base]{\color{textcolor}{\rmfamily\fontsize{8.000000}{9.600000}\selectfont\catcode`\^=\active\def^{\ifmmode\sp\else\^{}\fi}\catcode`\%=\active\def%{\%}$\mathdefault{2.97}$}}%
\end{pgfscope}%
\begin{pgfscope}%
\pgfpathrectangle{\pgfqpoint{0.724619in}{1.415315in}}{\pgfqpoint{2.720000in}{0.807339in}}%
\pgfusepath{clip}%
\pgfsetbuttcap%
\pgfsetroundjoin%
\pgfsetlinewidth{0.501875pt}%
\definecolor{currentstroke}{rgb}{0.690196,0.690196,0.690196}%
\pgfsetstrokecolor{currentstroke}%
\pgfsetstrokeopacity{0.250000}%
\pgfsetdash{{1.850000pt}{0.800000pt}}{0.000000pt}%
\pgfpathmoveto{\pgfqpoint{0.724619in}{1.729396in}}%
\pgfpathlineto{\pgfqpoint{3.444619in}{1.729396in}}%
\pgfusepath{stroke}%
\end{pgfscope}%
\begin{pgfscope}%
\pgfsetbuttcap%
\pgfsetroundjoin%
\definecolor{currentfill}{rgb}{0.000000,0.000000,0.000000}%
\pgfsetfillcolor{currentfill}%
\pgfsetlinewidth{0.803000pt}%
\definecolor{currentstroke}{rgb}{0.000000,0.000000,0.000000}%
\pgfsetstrokecolor{currentstroke}%
\pgfsetdash{}{0pt}%
\pgfsys@defobject{currentmarker}{\pgfqpoint{-0.048611in}{0.000000in}}{\pgfqpoint{-0.000000in}{0.000000in}}{%
\pgfpathmoveto{\pgfqpoint{-0.000000in}{0.000000in}}%
\pgfpathlineto{\pgfqpoint{-0.048611in}{0.000000in}}%
\pgfusepath{stroke,fill}%
}%
\begin{pgfscope}%
\pgfsys@transformshift{0.724619in}{1.729396in}%
\pgfsys@useobject{currentmarker}{}%
\end{pgfscope}%
\end{pgfscope}%
\begin{pgfscope}%
\definecolor{textcolor}{rgb}{0.000000,0.000000,0.000000}%
\pgfsetstrokecolor{textcolor}%
\pgfsetfillcolor{textcolor}%
\pgftext[x=0.417518in, y=1.690816in, left, base]{\color{textcolor}{\rmfamily\fontsize{8.000000}{9.600000}\selectfont\catcode`\^=\active\def^{\ifmmode\sp\else\^{}\fi}\catcode`\%=\active\def%{\%}$\mathdefault{2.98}$}}%
\end{pgfscope}%
\begin{pgfscope}%
\pgfpathrectangle{\pgfqpoint{0.724619in}{1.415315in}}{\pgfqpoint{2.720000in}{0.807339in}}%
\pgfusepath{clip}%
\pgfsetbuttcap%
\pgfsetroundjoin%
\pgfsetlinewidth{0.501875pt}%
\definecolor{currentstroke}{rgb}{0.690196,0.690196,0.690196}%
\pgfsetstrokecolor{currentstroke}%
\pgfsetstrokeopacity{0.250000}%
\pgfsetdash{{1.850000pt}{0.800000pt}}{0.000000pt}%
\pgfpathmoveto{\pgfqpoint{0.724619in}{1.957690in}}%
\pgfpathlineto{\pgfqpoint{3.444619in}{1.957690in}}%
\pgfusepath{stroke}%
\end{pgfscope}%
\begin{pgfscope}%
\pgfsetbuttcap%
\pgfsetroundjoin%
\definecolor{currentfill}{rgb}{0.000000,0.000000,0.000000}%
\pgfsetfillcolor{currentfill}%
\pgfsetlinewidth{0.803000pt}%
\definecolor{currentstroke}{rgb}{0.000000,0.000000,0.000000}%
\pgfsetstrokecolor{currentstroke}%
\pgfsetdash{}{0pt}%
\pgfsys@defobject{currentmarker}{\pgfqpoint{-0.048611in}{0.000000in}}{\pgfqpoint{-0.000000in}{0.000000in}}{%
\pgfpathmoveto{\pgfqpoint{-0.000000in}{0.000000in}}%
\pgfpathlineto{\pgfqpoint{-0.048611in}{0.000000in}}%
\pgfusepath{stroke,fill}%
}%
\begin{pgfscope}%
\pgfsys@transformshift{0.724619in}{1.957690in}%
\pgfsys@useobject{currentmarker}{}%
\end{pgfscope}%
\end{pgfscope}%
\begin{pgfscope}%
\definecolor{textcolor}{rgb}{0.000000,0.000000,0.000000}%
\pgfsetstrokecolor{textcolor}%
\pgfsetfillcolor{textcolor}%
\pgftext[x=0.417518in, y=1.919109in, left, base]{\color{textcolor}{\rmfamily\fontsize{8.000000}{9.600000}\selectfont\catcode`\^=\active\def^{\ifmmode\sp\else\^{}\fi}\catcode`\%=\active\def%{\%}$\mathdefault{2.99}$}}%
\end{pgfscope}%
\begin{pgfscope}%
\pgfpathrectangle{\pgfqpoint{0.724619in}{1.415315in}}{\pgfqpoint{2.720000in}{0.807339in}}%
\pgfusepath{clip}%
\pgfsetbuttcap%
\pgfsetroundjoin%
\pgfsetlinewidth{0.501875pt}%
\definecolor{currentstroke}{rgb}{0.690196,0.690196,0.690196}%
\pgfsetstrokecolor{currentstroke}%
\pgfsetstrokeopacity{0.250000}%
\pgfsetdash{{1.850000pt}{0.800000pt}}{0.000000pt}%
\pgfpathmoveto{\pgfqpoint{0.724619in}{2.185983in}}%
\pgfpathlineto{\pgfqpoint{3.444619in}{2.185983in}}%
\pgfusepath{stroke}%
\end{pgfscope}%
\begin{pgfscope}%
\pgfsetbuttcap%
\pgfsetroundjoin%
\definecolor{currentfill}{rgb}{0.000000,0.000000,0.000000}%
\pgfsetfillcolor{currentfill}%
\pgfsetlinewidth{0.803000pt}%
\definecolor{currentstroke}{rgb}{0.000000,0.000000,0.000000}%
\pgfsetstrokecolor{currentstroke}%
\pgfsetdash{}{0pt}%
\pgfsys@defobject{currentmarker}{\pgfqpoint{-0.048611in}{0.000000in}}{\pgfqpoint{-0.000000in}{0.000000in}}{%
\pgfpathmoveto{\pgfqpoint{-0.000000in}{0.000000in}}%
\pgfpathlineto{\pgfqpoint{-0.048611in}{0.000000in}}%
\pgfusepath{stroke,fill}%
}%
\begin{pgfscope}%
\pgfsys@transformshift{0.724619in}{2.185983in}%
\pgfsys@useobject{currentmarker}{}%
\end{pgfscope}%
\end{pgfscope}%
\begin{pgfscope}%
\definecolor{textcolor}{rgb}{0.000000,0.000000,0.000000}%
\pgfsetstrokecolor{textcolor}%
\pgfsetfillcolor{textcolor}%
\pgftext[x=0.417518in, y=2.147403in, left, base]{\color{textcolor}{\rmfamily\fontsize{8.000000}{9.600000}\selectfont\catcode`\^=\active\def^{\ifmmode\sp\else\^{}\fi}\catcode`\%=\active\def%{\%}$\mathdefault{3.00}$}}%
\end{pgfscope}%
\begin{pgfscope}%
\definecolor{textcolor}{rgb}{0.000000,0.000000,0.000000}%
\pgfsetstrokecolor{textcolor}%
\pgfsetfillcolor{textcolor}%
\pgftext[x=0.361962in,y=1.818985in,,bottom,rotate=90.000000]{\color{textcolor}{\rmfamily\fontsize{8.000000}{9.600000}\selectfont\catcode`\^=\active\def^{\ifmmode\sp\else\^{}\fi}\catcode`\%=\active\def%{\%}$i_{\mathrm{in}}$ (A)}}%
\end{pgfscope}%
\begin{pgfscope}%
\pgfpathrectangle{\pgfqpoint{0.724619in}{1.415315in}}{\pgfqpoint{2.720000in}{0.807339in}}%
\pgfusepath{clip}%
\pgfsetrectcap%
\pgfsetroundjoin%
\pgfsetlinewidth{1.505625pt}%
\definecolor{currentstroke}{rgb}{0.835294,0.368627,0.000000}%
\pgfsetstrokecolor{currentstroke}%
\pgfsetdash{}{0pt}%
\pgfpathmoveto{\pgfqpoint{0.848256in}{1.452055in}}%
\pgfpathlineto{\pgfqpoint{0.848266in}{1.452019in}}%
\pgfpathlineto{\pgfqpoint{0.848543in}{1.453154in}}%
\pgfpathlineto{\pgfqpoint{1.024931in}{2.185957in}}%
\pgfpathlineto{\pgfqpoint{1.025111in}{2.185396in}}%
\pgfpathlineto{\pgfqpoint{1.244647in}{1.501253in}}%
\pgfpathlineto{\pgfqpoint{1.260388in}{1.452017in}}%
\pgfpathlineto{\pgfqpoint{1.260665in}{1.453153in}}%
\pgfpathlineto{\pgfqpoint{1.437052in}{2.185956in}}%
\pgfpathlineto{\pgfqpoint{1.437232in}{2.185394in}}%
\pgfpathlineto{\pgfqpoint{1.656768in}{1.501251in}}%
\pgfpathlineto{\pgfqpoint{1.672509in}{1.452016in}}%
\pgfpathlineto{\pgfqpoint{1.672786in}{1.453151in}}%
\pgfpathlineto{\pgfqpoint{1.849174in}{2.185954in}}%
\pgfpathlineto{\pgfqpoint{1.849353in}{2.185393in}}%
\pgfpathlineto{\pgfqpoint{2.084630in}{1.452015in}}%
\pgfpathlineto{\pgfqpoint{2.085168in}{1.454237in}}%
\pgfpathlineto{\pgfqpoint{2.261295in}{2.185953in}}%
\pgfpathlineto{\pgfqpoint{2.261474in}{2.185392in}}%
\pgfpathlineto{\pgfqpoint{2.496751in}{1.452013in}}%
\pgfpathlineto{\pgfqpoint{2.497289in}{1.454236in}}%
\pgfpathlineto{\pgfqpoint{2.673416in}{2.185951in}}%
\pgfpathlineto{\pgfqpoint{2.673595in}{2.185390in}}%
\pgfpathlineto{\pgfqpoint{2.893131in}{1.501247in}}%
\pgfpathlineto{\pgfqpoint{2.908872in}{1.452012in}}%
\pgfpathlineto{\pgfqpoint{2.909149in}{1.453147in}}%
\pgfpathlineto{\pgfqpoint{3.085537in}{2.185950in}}%
\pgfpathlineto{\pgfqpoint{3.085717in}{2.185389in}}%
\pgfpathlineto{\pgfqpoint{3.305253in}{1.501246in}}%
\pgfpathlineto{\pgfqpoint{3.320983in}{1.452047in}}%
\pgfpathlineto{\pgfqpoint{3.320983in}{1.452047in}}%
\pgfusepath{stroke}%
\end{pgfscope}%
\begin{pgfscope}%
\pgfsetrectcap%
\pgfsetmiterjoin%
\pgfsetlinewidth{0.803000pt}%
\definecolor{currentstroke}{rgb}{0.000000,0.000000,0.000000}%
\pgfsetstrokecolor{currentstroke}%
\pgfsetdash{}{0pt}%
\pgfpathmoveto{\pgfqpoint{0.724619in}{1.415315in}}%
\pgfpathlineto{\pgfqpoint{0.724619in}{2.222654in}}%
\pgfusepath{stroke}%
\end{pgfscope}%
\begin{pgfscope}%
\pgfsetrectcap%
\pgfsetmiterjoin%
\pgfsetlinewidth{0.803000pt}%
\definecolor{currentstroke}{rgb}{0.000000,0.000000,0.000000}%
\pgfsetstrokecolor{currentstroke}%
\pgfsetdash{}{0pt}%
\pgfpathmoveto{\pgfqpoint{3.444619in}{1.415315in}}%
\pgfpathlineto{\pgfqpoint{3.444619in}{2.222654in}}%
\pgfusepath{stroke}%
\end{pgfscope}%
\begin{pgfscope}%
\pgfsetrectcap%
\pgfsetmiterjoin%
\pgfsetlinewidth{0.803000pt}%
\definecolor{currentstroke}{rgb}{0.000000,0.000000,0.000000}%
\pgfsetstrokecolor{currentstroke}%
\pgfsetdash{}{0pt}%
\pgfpathmoveto{\pgfqpoint{0.724619in}{1.415315in}}%
\pgfpathlineto{\pgfqpoint{3.444619in}{1.415315in}}%
\pgfusepath{stroke}%
\end{pgfscope}%
\begin{pgfscope}%
\pgfsetrectcap%
\pgfsetmiterjoin%
\pgfsetlinewidth{0.803000pt}%
\definecolor{currentstroke}{rgb}{0.000000,0.000000,0.000000}%
\pgfsetstrokecolor{currentstroke}%
\pgfsetdash{}{0pt}%
\pgfpathmoveto{\pgfqpoint{0.724619in}{2.222654in}}%
\pgfpathlineto{\pgfqpoint{3.444619in}{2.222654in}}%
\pgfusepath{stroke}%
\end{pgfscope}%
\begin{pgfscope}%
\definecolor{textcolor}{rgb}{0.000000,0.000000,0.000000}%
\pgfsetstrokecolor{textcolor}%
\pgfsetfillcolor{textcolor}%
\pgftext[x=3.254219in,y=2.166141in,left,top]{\color{textcolor}{\rmfamily\fontsize{8.000000}{9.600000}\selectfont\catcode`\^=\active\def^{\ifmmode\sp\else\^{}\fi}\catcode`\%=\active\def%{\%}(a)}}%
\end{pgfscope}%
\begin{pgfscope}%
\pgfsetbuttcap%
\pgfsetmiterjoin%
\definecolor{currentfill}{rgb}{1.000000,1.000000,1.000000}%
\pgfsetfillcolor{currentfill}%
\pgfsetlinewidth{0.000000pt}%
\definecolor{currentstroke}{rgb}{0.000000,0.000000,0.000000}%
\pgfsetstrokecolor{currentstroke}%
\pgfsetstrokeopacity{0.000000}%
\pgfsetdash{}{0pt}%
\pgfpathmoveto{\pgfqpoint{0.724619in}{0.462654in}}%
\pgfpathlineto{\pgfqpoint{3.444619in}{0.462654in}}%
\pgfpathlineto{\pgfqpoint{3.444619in}{1.269994in}}%
\pgfpathlineto{\pgfqpoint{0.724619in}{1.269994in}}%
\pgfpathlineto{\pgfqpoint{0.724619in}{0.462654in}}%
\pgfpathclose%
\pgfusepath{fill}%
\end{pgfscope}%
\begin{pgfscope}%
\pgfpathrectangle{\pgfqpoint{0.724619in}{0.462654in}}{\pgfqpoint{2.720000in}{0.807339in}}%
\pgfusepath{clip}%
\pgfsetbuttcap%
\pgfsetroundjoin%
\pgfsetlinewidth{0.501875pt}%
\definecolor{currentstroke}{rgb}{0.690196,0.690196,0.690196}%
\pgfsetstrokecolor{currentstroke}%
\pgfsetstrokeopacity{0.250000}%
\pgfsetdash{{1.850000pt}{0.800000pt}}{0.000000pt}%
\pgfpathmoveto{\pgfqpoint{1.079044in}{0.462654in}}%
\pgfpathlineto{\pgfqpoint{1.079044in}{1.269994in}}%
\pgfusepath{stroke}%
\end{pgfscope}%
\begin{pgfscope}%
\pgfsetbuttcap%
\pgfsetroundjoin%
\definecolor{currentfill}{rgb}{0.000000,0.000000,0.000000}%
\pgfsetfillcolor{currentfill}%
\pgfsetlinewidth{0.803000pt}%
\definecolor{currentstroke}{rgb}{0.000000,0.000000,0.000000}%
\pgfsetstrokecolor{currentstroke}%
\pgfsetdash{}{0pt}%
\pgfsys@defobject{currentmarker}{\pgfqpoint{0.000000in}{-0.048611in}}{\pgfqpoint{0.000000in}{0.000000in}}{%
\pgfpathmoveto{\pgfqpoint{0.000000in}{0.000000in}}%
\pgfpathlineto{\pgfqpoint{0.000000in}{-0.048611in}}%
\pgfusepath{stroke,fill}%
}%
\begin{pgfscope}%
\pgfsys@transformshift{1.079044in}{0.462654in}%
\pgfsys@useobject{currentmarker}{}%
\end{pgfscope}%
\end{pgfscope}%
\begin{pgfscope}%
\definecolor{textcolor}{rgb}{0.000000,0.000000,0.000000}%
\pgfsetstrokecolor{textcolor}%
\pgfsetfillcolor{textcolor}%
\pgftext[x=1.079044in,y=0.365432in,,top]{\color{textcolor}{\rmfamily\fontsize{8.000000}{9.600000}\selectfont\catcode`\^=\active\def^{\ifmmode\sp\else\^{}\fi}\catcode`\%=\active\def%{\%}$\mathdefault{78.04}$}}%
\end{pgfscope}%
\begin{pgfscope}%
\pgfpathrectangle{\pgfqpoint{0.724619in}{0.462654in}}{\pgfqpoint{2.720000in}{0.807339in}}%
\pgfusepath{clip}%
\pgfsetbuttcap%
\pgfsetroundjoin%
\pgfsetlinewidth{0.501875pt}%
\definecolor{currentstroke}{rgb}{0.690196,0.690196,0.690196}%
\pgfsetstrokecolor{currentstroke}%
\pgfsetstrokeopacity{0.250000}%
\pgfsetdash{{1.850000pt}{0.800000pt}}{0.000000pt}%
\pgfpathmoveto{\pgfqpoint{1.606559in}{0.462654in}}%
\pgfpathlineto{\pgfqpoint{1.606559in}{1.269994in}}%
\pgfusepath{stroke}%
\end{pgfscope}%
\begin{pgfscope}%
\pgfsetbuttcap%
\pgfsetroundjoin%
\definecolor{currentfill}{rgb}{0.000000,0.000000,0.000000}%
\pgfsetfillcolor{currentfill}%
\pgfsetlinewidth{0.803000pt}%
\definecolor{currentstroke}{rgb}{0.000000,0.000000,0.000000}%
\pgfsetstrokecolor{currentstroke}%
\pgfsetdash{}{0pt}%
\pgfsys@defobject{currentmarker}{\pgfqpoint{0.000000in}{-0.048611in}}{\pgfqpoint{0.000000in}{0.000000in}}{%
\pgfpathmoveto{\pgfqpoint{0.000000in}{0.000000in}}%
\pgfpathlineto{\pgfqpoint{0.000000in}{-0.048611in}}%
\pgfusepath{stroke,fill}%
}%
\begin{pgfscope}%
\pgfsys@transformshift{1.606559in}{0.462654in}%
\pgfsys@useobject{currentmarker}{}%
\end{pgfscope}%
\end{pgfscope}%
\begin{pgfscope}%
\definecolor{textcolor}{rgb}{0.000000,0.000000,0.000000}%
\pgfsetstrokecolor{textcolor}%
\pgfsetfillcolor{textcolor}%
\pgftext[x=1.606559in,y=0.365432in,,top]{\color{textcolor}{\rmfamily\fontsize{8.000000}{9.600000}\selectfont\catcode`\^=\active\def^{\ifmmode\sp\else\^{}\fi}\catcode`\%=\active\def%{\%}$\mathdefault{78.06}$}}%
\end{pgfscope}%
\begin{pgfscope}%
\pgfpathrectangle{\pgfqpoint{0.724619in}{0.462654in}}{\pgfqpoint{2.720000in}{0.807339in}}%
\pgfusepath{clip}%
\pgfsetbuttcap%
\pgfsetroundjoin%
\pgfsetlinewidth{0.501875pt}%
\definecolor{currentstroke}{rgb}{0.690196,0.690196,0.690196}%
\pgfsetstrokecolor{currentstroke}%
\pgfsetstrokeopacity{0.250000}%
\pgfsetdash{{1.850000pt}{0.800000pt}}{0.000000pt}%
\pgfpathmoveto{\pgfqpoint{2.134074in}{0.462654in}}%
\pgfpathlineto{\pgfqpoint{2.134074in}{1.269994in}}%
\pgfusepath{stroke}%
\end{pgfscope}%
\begin{pgfscope}%
\pgfsetbuttcap%
\pgfsetroundjoin%
\definecolor{currentfill}{rgb}{0.000000,0.000000,0.000000}%
\pgfsetfillcolor{currentfill}%
\pgfsetlinewidth{0.803000pt}%
\definecolor{currentstroke}{rgb}{0.000000,0.000000,0.000000}%
\pgfsetstrokecolor{currentstroke}%
\pgfsetdash{}{0pt}%
\pgfsys@defobject{currentmarker}{\pgfqpoint{0.000000in}{-0.048611in}}{\pgfqpoint{0.000000in}{0.000000in}}{%
\pgfpathmoveto{\pgfqpoint{0.000000in}{0.000000in}}%
\pgfpathlineto{\pgfqpoint{0.000000in}{-0.048611in}}%
\pgfusepath{stroke,fill}%
}%
\begin{pgfscope}%
\pgfsys@transformshift{2.134074in}{0.462654in}%
\pgfsys@useobject{currentmarker}{}%
\end{pgfscope}%
\end{pgfscope}%
\begin{pgfscope}%
\definecolor{textcolor}{rgb}{0.000000,0.000000,0.000000}%
\pgfsetstrokecolor{textcolor}%
\pgfsetfillcolor{textcolor}%
\pgftext[x=2.134074in,y=0.365432in,,top]{\color{textcolor}{\rmfamily\fontsize{8.000000}{9.600000}\selectfont\catcode`\^=\active\def^{\ifmmode\sp\else\^{}\fi}\catcode`\%=\active\def%{\%}$\mathdefault{78.08}$}}%
\end{pgfscope}%
\begin{pgfscope}%
\pgfpathrectangle{\pgfqpoint{0.724619in}{0.462654in}}{\pgfqpoint{2.720000in}{0.807339in}}%
\pgfusepath{clip}%
\pgfsetbuttcap%
\pgfsetroundjoin%
\pgfsetlinewidth{0.501875pt}%
\definecolor{currentstroke}{rgb}{0.690196,0.690196,0.690196}%
\pgfsetstrokecolor{currentstroke}%
\pgfsetstrokeopacity{0.250000}%
\pgfsetdash{{1.850000pt}{0.800000pt}}{0.000000pt}%
\pgfpathmoveto{\pgfqpoint{2.661589in}{0.462654in}}%
\pgfpathlineto{\pgfqpoint{2.661589in}{1.269994in}}%
\pgfusepath{stroke}%
\end{pgfscope}%
\begin{pgfscope}%
\pgfsetbuttcap%
\pgfsetroundjoin%
\definecolor{currentfill}{rgb}{0.000000,0.000000,0.000000}%
\pgfsetfillcolor{currentfill}%
\pgfsetlinewidth{0.803000pt}%
\definecolor{currentstroke}{rgb}{0.000000,0.000000,0.000000}%
\pgfsetstrokecolor{currentstroke}%
\pgfsetdash{}{0pt}%
\pgfsys@defobject{currentmarker}{\pgfqpoint{0.000000in}{-0.048611in}}{\pgfqpoint{0.000000in}{0.000000in}}{%
\pgfpathmoveto{\pgfqpoint{0.000000in}{0.000000in}}%
\pgfpathlineto{\pgfqpoint{0.000000in}{-0.048611in}}%
\pgfusepath{stroke,fill}%
}%
\begin{pgfscope}%
\pgfsys@transformshift{2.661589in}{0.462654in}%
\pgfsys@useobject{currentmarker}{}%
\end{pgfscope}%
\end{pgfscope}%
\begin{pgfscope}%
\definecolor{textcolor}{rgb}{0.000000,0.000000,0.000000}%
\pgfsetstrokecolor{textcolor}%
\pgfsetfillcolor{textcolor}%
\pgftext[x=2.661589in,y=0.365432in,,top]{\color{textcolor}{\rmfamily\fontsize{8.000000}{9.600000}\selectfont\catcode`\^=\active\def^{\ifmmode\sp\else\^{}\fi}\catcode`\%=\active\def%{\%}$\mathdefault{78.10}$}}%
\end{pgfscope}%
\begin{pgfscope}%
\pgfpathrectangle{\pgfqpoint{0.724619in}{0.462654in}}{\pgfqpoint{2.720000in}{0.807339in}}%
\pgfusepath{clip}%
\pgfsetbuttcap%
\pgfsetroundjoin%
\pgfsetlinewidth{0.501875pt}%
\definecolor{currentstroke}{rgb}{0.690196,0.690196,0.690196}%
\pgfsetstrokecolor{currentstroke}%
\pgfsetstrokeopacity{0.250000}%
\pgfsetdash{{1.850000pt}{0.800000pt}}{0.000000pt}%
\pgfpathmoveto{\pgfqpoint{3.189104in}{0.462654in}}%
\pgfpathlineto{\pgfqpoint{3.189104in}{1.269994in}}%
\pgfusepath{stroke}%
\end{pgfscope}%
\begin{pgfscope}%
\pgfsetbuttcap%
\pgfsetroundjoin%
\definecolor{currentfill}{rgb}{0.000000,0.000000,0.000000}%
\pgfsetfillcolor{currentfill}%
\pgfsetlinewidth{0.803000pt}%
\definecolor{currentstroke}{rgb}{0.000000,0.000000,0.000000}%
\pgfsetstrokecolor{currentstroke}%
\pgfsetdash{}{0pt}%
\pgfsys@defobject{currentmarker}{\pgfqpoint{0.000000in}{-0.048611in}}{\pgfqpoint{0.000000in}{0.000000in}}{%
\pgfpathmoveto{\pgfqpoint{0.000000in}{0.000000in}}%
\pgfpathlineto{\pgfqpoint{0.000000in}{-0.048611in}}%
\pgfusepath{stroke,fill}%
}%
\begin{pgfscope}%
\pgfsys@transformshift{3.189104in}{0.462654in}%
\pgfsys@useobject{currentmarker}{}%
\end{pgfscope}%
\end{pgfscope}%
\begin{pgfscope}%
\definecolor{textcolor}{rgb}{0.000000,0.000000,0.000000}%
\pgfsetstrokecolor{textcolor}%
\pgfsetfillcolor{textcolor}%
\pgftext[x=3.189104in,y=0.365432in,,top]{\color{textcolor}{\rmfamily\fontsize{8.000000}{9.600000}\selectfont\catcode`\^=\active\def^{\ifmmode\sp\else\^{}\fi}\catcode`\%=\active\def%{\%}$\mathdefault{78.12}$}}%
\end{pgfscope}%
\begin{pgfscope}%
\definecolor{textcolor}{rgb}{0.000000,0.000000,0.000000}%
\pgfsetstrokecolor{textcolor}%
\pgfsetfillcolor{textcolor}%
\pgftext[x=2.084619in,y=0.211111in,,top]{\color{textcolor}{\rmfamily\fontsize{8.000000}{9.600000}\selectfont\catcode`\^=\active\def^{\ifmmode\sp\else\^{}\fi}\catcode`\%=\active\def%{\%}Time (ms)}}%
\end{pgfscope}%
\begin{pgfscope}%
\pgfpathrectangle{\pgfqpoint{0.724619in}{0.462654in}}{\pgfqpoint{2.720000in}{0.807339in}}%
\pgfusepath{clip}%
\pgfsetbuttcap%
\pgfsetroundjoin%
\pgfsetlinewidth{0.501875pt}%
\definecolor{currentstroke}{rgb}{0.690196,0.690196,0.690196}%
\pgfsetstrokecolor{currentstroke}%
\pgfsetstrokeopacity{0.250000}%
\pgfsetdash{{1.850000pt}{0.800000pt}}{0.000000pt}%
\pgfpathmoveto{\pgfqpoint{0.724619in}{0.528440in}}%
\pgfpathlineto{\pgfqpoint{3.444619in}{0.528440in}}%
\pgfusepath{stroke}%
\end{pgfscope}%
\begin{pgfscope}%
\pgfsetbuttcap%
\pgfsetroundjoin%
\definecolor{currentfill}{rgb}{0.000000,0.000000,0.000000}%
\pgfsetfillcolor{currentfill}%
\pgfsetlinewidth{0.803000pt}%
\definecolor{currentstroke}{rgb}{0.000000,0.000000,0.000000}%
\pgfsetstrokecolor{currentstroke}%
\pgfsetdash{}{0pt}%
\pgfsys@defobject{currentmarker}{\pgfqpoint{-0.048611in}{0.000000in}}{\pgfqpoint{-0.000000in}{0.000000in}}{%
\pgfpathmoveto{\pgfqpoint{-0.000000in}{0.000000in}}%
\pgfpathlineto{\pgfqpoint{-0.048611in}{0.000000in}}%
\pgfusepath{stroke,fill}%
}%
\begin{pgfscope}%
\pgfsys@transformshift{0.724619in}{0.528440in}%
\pgfsys@useobject{currentmarker}{}%
\end{pgfscope}%
\end{pgfscope}%
\begin{pgfscope}%
\definecolor{textcolor}{rgb}{0.000000,0.000000,0.000000}%
\pgfsetstrokecolor{textcolor}%
\pgfsetfillcolor{textcolor}%
\pgftext[x=0.266667in, y=0.489860in, left, base]{\color{textcolor}{\rmfamily\fontsize{8.000000}{9.600000}\selectfont\catcode`\^=\active\def^{\ifmmode\sp\else\^{}\fi}\catcode`\%=\active\def%{\%}$\mathdefault{\ensuremath{-}4.000}$}}%
\end{pgfscope}%
\begin{pgfscope}%
\pgfpathrectangle{\pgfqpoint{0.724619in}{0.462654in}}{\pgfqpoint{2.720000in}{0.807339in}}%
\pgfusepath{clip}%
\pgfsetbuttcap%
\pgfsetroundjoin%
\pgfsetlinewidth{0.501875pt}%
\definecolor{currentstroke}{rgb}{0.690196,0.690196,0.690196}%
\pgfsetstrokecolor{currentstroke}%
\pgfsetstrokeopacity{0.250000}%
\pgfsetdash{{1.850000pt}{0.800000pt}}{0.000000pt}%
\pgfpathmoveto{\pgfqpoint{0.724619in}{0.863731in}}%
\pgfpathlineto{\pgfqpoint{3.444619in}{0.863731in}}%
\pgfusepath{stroke}%
\end{pgfscope}%
\begin{pgfscope}%
\pgfsetbuttcap%
\pgfsetroundjoin%
\definecolor{currentfill}{rgb}{0.000000,0.000000,0.000000}%
\pgfsetfillcolor{currentfill}%
\pgfsetlinewidth{0.803000pt}%
\definecolor{currentstroke}{rgb}{0.000000,0.000000,0.000000}%
\pgfsetstrokecolor{currentstroke}%
\pgfsetdash{}{0pt}%
\pgfsys@defobject{currentmarker}{\pgfqpoint{-0.048611in}{0.000000in}}{\pgfqpoint{-0.000000in}{0.000000in}}{%
\pgfpathmoveto{\pgfqpoint{-0.000000in}{0.000000in}}%
\pgfpathlineto{\pgfqpoint{-0.048611in}{0.000000in}}%
\pgfusepath{stroke,fill}%
}%
\begin{pgfscope}%
\pgfsys@transformshift{0.724619in}{0.863731in}%
\pgfsys@useobject{currentmarker}{}%
\end{pgfscope}%
\end{pgfscope}%
\begin{pgfscope}%
\definecolor{textcolor}{rgb}{0.000000,0.000000,0.000000}%
\pgfsetstrokecolor{textcolor}%
\pgfsetfillcolor{textcolor}%
\pgftext[x=0.266667in, y=0.825151in, left, base]{\color{textcolor}{\rmfamily\fontsize{8.000000}{9.600000}\selectfont\catcode`\^=\active\def^{\ifmmode\sp\else\^{}\fi}\catcode`\%=\active\def%{\%}$\mathdefault{\ensuremath{-}3.975}$}}%
\end{pgfscope}%
\begin{pgfscope}%
\pgfpathrectangle{\pgfqpoint{0.724619in}{0.462654in}}{\pgfqpoint{2.720000in}{0.807339in}}%
\pgfusepath{clip}%
\pgfsetbuttcap%
\pgfsetroundjoin%
\pgfsetlinewidth{0.501875pt}%
\definecolor{currentstroke}{rgb}{0.690196,0.690196,0.690196}%
\pgfsetstrokecolor{currentstroke}%
\pgfsetstrokeopacity{0.250000}%
\pgfsetdash{{1.850000pt}{0.800000pt}}{0.000000pt}%
\pgfpathmoveto{\pgfqpoint{0.724619in}{1.199023in}}%
\pgfpathlineto{\pgfqpoint{3.444619in}{1.199023in}}%
\pgfusepath{stroke}%
\end{pgfscope}%
\begin{pgfscope}%
\pgfsetbuttcap%
\pgfsetroundjoin%
\definecolor{currentfill}{rgb}{0.000000,0.000000,0.000000}%
\pgfsetfillcolor{currentfill}%
\pgfsetlinewidth{0.803000pt}%
\definecolor{currentstroke}{rgb}{0.000000,0.000000,0.000000}%
\pgfsetstrokecolor{currentstroke}%
\pgfsetdash{}{0pt}%
\pgfsys@defobject{currentmarker}{\pgfqpoint{-0.048611in}{0.000000in}}{\pgfqpoint{-0.000000in}{0.000000in}}{%
\pgfpathmoveto{\pgfqpoint{-0.000000in}{0.000000in}}%
\pgfpathlineto{\pgfqpoint{-0.048611in}{0.000000in}}%
\pgfusepath{stroke,fill}%
}%
\begin{pgfscope}%
\pgfsys@transformshift{0.724619in}{1.199023in}%
\pgfsys@useobject{currentmarker}{}%
\end{pgfscope}%
\end{pgfscope}%
\begin{pgfscope}%
\definecolor{textcolor}{rgb}{0.000000,0.000000,0.000000}%
\pgfsetstrokecolor{textcolor}%
\pgfsetfillcolor{textcolor}%
\pgftext[x=0.266667in, y=1.160443in, left, base]{\color{textcolor}{\rmfamily\fontsize{8.000000}{9.600000}\selectfont\catcode`\^=\active\def^{\ifmmode\sp\else\^{}\fi}\catcode`\%=\active\def%{\%}$\mathdefault{\ensuremath{-}3.950}$}}%
\end{pgfscope}%
\begin{pgfscope}%
\definecolor{textcolor}{rgb}{0.000000,0.000000,0.000000}%
\pgfsetstrokecolor{textcolor}%
\pgfsetfillcolor{textcolor}%
\pgftext[x=0.211111in,y=0.866324in,,bottom,rotate=90.000000]{\color{textcolor}{\rmfamily\fontsize{8.000000}{9.600000}\selectfont\catcode`\^=\active\def^{\ifmmode\sp\else\^{}\fi}\catcode`\%=\active\def%{\%}$i_{\mathrm{out}}$ (A)}}%
\end{pgfscope}%
\begin{pgfscope}%
\pgfpathrectangle{\pgfqpoint{0.724619in}{0.462654in}}{\pgfqpoint{2.720000in}{0.807339in}}%
\pgfusepath{clip}%
\pgfsetrectcap%
\pgfsetroundjoin%
\pgfsetlinewidth{1.505625pt}%
\definecolor{currentstroke}{rgb}{0.835294,0.368627,0.000000}%
\pgfsetstrokecolor{currentstroke}%
\pgfsetdash{}{0pt}%
\pgfpathmoveto{\pgfqpoint{0.848256in}{0.998027in}}%
\pgfpathlineto{\pgfqpoint{0.856588in}{1.050769in}}%
\pgfpathlineto{\pgfqpoint{0.864830in}{1.095966in}}%
\pgfpathlineto{\pgfqpoint{0.873073in}{1.134413in}}%
\pgfpathlineto{\pgfqpoint{0.881315in}{1.166294in}}%
\pgfpathlineto{\pgfqpoint{0.889558in}{1.191786in}}%
\pgfpathlineto{\pgfqpoint{0.897800in}{1.211064in}}%
\pgfpathlineto{\pgfqpoint{0.906042in}{1.224296in}}%
\pgfpathlineto{\pgfqpoint{0.910165in}{1.228696in}}%
\pgfpathlineto{\pgfqpoint{0.914285in}{1.231647in}}%
\pgfpathlineto{\pgfqpoint{0.918407in}{1.233168in}}%
\pgfpathlineto{\pgfqpoint{0.922527in}{1.233280in}}%
\pgfpathlineto{\pgfqpoint{0.926650in}{1.232000in}}%
\pgfpathlineto{\pgfqpoint{0.930770in}{1.229350in}}%
\pgfpathlineto{\pgfqpoint{0.934892in}{1.225348in}}%
\pgfpathlineto{\pgfqpoint{0.939012in}{1.220012in}}%
\pgfpathlineto{\pgfqpoint{0.947255in}{1.205415in}}%
\pgfpathlineto{\pgfqpoint{0.955497in}{1.185705in}}%
\pgfpathlineto{\pgfqpoint{0.963739in}{1.161024in}}%
\pgfpathlineto{\pgfqpoint{0.971982in}{1.131511in}}%
\pgfpathlineto{\pgfqpoint{0.980224in}{1.097302in}}%
\pgfpathlineto{\pgfqpoint{0.992589in}{1.037471in}}%
\pgfpathlineto{\pgfqpoint{1.004952in}{0.967803in}}%
\pgfpathlineto{\pgfqpoint{1.017317in}{0.888712in}}%
\pgfpathlineto{\pgfqpoint{1.038586in}{0.744887in}}%
\pgfpathlineto{\pgfqpoint{1.050951in}{0.676051in}}%
\pgfpathlineto{\pgfqpoint{1.059193in}{0.636869in}}%
\pgfpathlineto{\pgfqpoint{1.067436in}{0.602888in}}%
\pgfpathlineto{\pgfqpoint{1.075678in}{0.573967in}}%
\pgfpathlineto{\pgfqpoint{1.083921in}{0.549972in}}%
\pgfpathlineto{\pgfqpoint{1.092163in}{0.530770in}}%
\pgfpathlineto{\pgfqpoint{1.100405in}{0.516231in}}%
\pgfpathlineto{\pgfqpoint{1.108648in}{0.506230in}}%
\pgfpathlineto{\pgfqpoint{1.112768in}{0.502893in}}%
\pgfpathlineto{\pgfqpoint{1.116890in}{0.500644in}}%
\pgfpathlineto{\pgfqpoint{1.121010in}{0.499468in}}%
\pgfpathlineto{\pgfqpoint{1.125133in}{0.499352in}}%
\pgfpathlineto{\pgfqpoint{1.129253in}{0.500280in}}%
\pgfpathlineto{\pgfqpoint{1.133375in}{0.502238in}}%
\pgfpathlineto{\pgfqpoint{1.137495in}{0.505211in}}%
\pgfpathlineto{\pgfqpoint{1.145737in}{0.514150in}}%
\pgfpathlineto{\pgfqpoint{1.153980in}{0.526987in}}%
\pgfpathlineto{\pgfqpoint{1.162222in}{0.543614in}}%
\pgfpathlineto{\pgfqpoint{1.170465in}{0.563927in}}%
\pgfpathlineto{\pgfqpoint{1.178707in}{0.587823in}}%
\pgfpathlineto{\pgfqpoint{1.191072in}{0.630169in}}%
\pgfpathlineto{\pgfqpoint{1.203434in}{0.680029in}}%
\pgfpathlineto{\pgfqpoint{1.215799in}{0.737090in}}%
\pgfpathlineto{\pgfqpoint{1.228162in}{0.801048in}}%
\pgfpathlineto{\pgfqpoint{1.244647in}{0.896554in}}%
\pgfpathlineto{\pgfqpoint{1.262495in}{1.012122in}}%
\pgfpathlineto{\pgfqpoint{1.272832in}{1.074226in}}%
\pgfpathlineto{\pgfqpoint{1.281074in}{1.116025in}}%
\pgfpathlineto{\pgfqpoint{1.289317in}{1.151167in}}%
\pgfpathlineto{\pgfqpoint{1.297559in}{1.179831in}}%
\pgfpathlineto{\pgfqpoint{1.305801in}{1.202195in}}%
\pgfpathlineto{\pgfqpoint{1.314044in}{1.218428in}}%
\pgfpathlineto{\pgfqpoint{1.318164in}{1.224299in}}%
\pgfpathlineto{\pgfqpoint{1.322286in}{1.228700in}}%
\pgfpathlineto{\pgfqpoint{1.326406in}{1.231651in}}%
\pgfpathlineto{\pgfqpoint{1.330529in}{1.233172in}}%
\pgfpathlineto{\pgfqpoint{1.334649in}{1.233283in}}%
\pgfpathlineto{\pgfqpoint{1.338771in}{1.232004in}}%
\pgfpathlineto{\pgfqpoint{1.342891in}{1.229353in}}%
\pgfpathlineto{\pgfqpoint{1.347013in}{1.225351in}}%
\pgfpathlineto{\pgfqpoint{1.351133in}{1.220015in}}%
\pgfpathlineto{\pgfqpoint{1.359376in}{1.205418in}}%
\pgfpathlineto{\pgfqpoint{1.367618in}{1.185708in}}%
\pgfpathlineto{\pgfqpoint{1.375861in}{1.161027in}}%
\pgfpathlineto{\pgfqpoint{1.384103in}{1.131515in}}%
\pgfpathlineto{\pgfqpoint{1.392345in}{1.097306in}}%
\pgfpathlineto{\pgfqpoint{1.404710in}{1.037475in}}%
\pgfpathlineto{\pgfqpoint{1.417073in}{0.967806in}}%
\pgfpathlineto{\pgfqpoint{1.429438in}{0.888716in}}%
\pgfpathlineto{\pgfqpoint{1.450707in}{0.744890in}}%
\pgfpathlineto{\pgfqpoint{1.463072in}{0.676054in}}%
\pgfpathlineto{\pgfqpoint{1.471315in}{0.636872in}}%
\pgfpathlineto{\pgfqpoint{1.479557in}{0.602891in}}%
\pgfpathlineto{\pgfqpoint{1.487799in}{0.573971in}}%
\pgfpathlineto{\pgfqpoint{1.496042in}{0.549975in}}%
\pgfpathlineto{\pgfqpoint{1.504284in}{0.530773in}}%
\pgfpathlineto{\pgfqpoint{1.512527in}{0.516235in}}%
\pgfpathlineto{\pgfqpoint{1.520769in}{0.506233in}}%
\pgfpathlineto{\pgfqpoint{1.524889in}{0.502896in}}%
\pgfpathlineto{\pgfqpoint{1.529011in}{0.500647in}}%
\pgfpathlineto{\pgfqpoint{1.533131in}{0.499472in}}%
\pgfpathlineto{\pgfqpoint{1.537254in}{0.499355in}}%
\pgfpathlineto{\pgfqpoint{1.541374in}{0.500283in}}%
\pgfpathlineto{\pgfqpoint{1.545496in}{0.502241in}}%
\pgfpathlineto{\pgfqpoint{1.549616in}{0.505215in}}%
\pgfpathlineto{\pgfqpoint{1.557859in}{0.514154in}}%
\pgfpathlineto{\pgfqpoint{1.566101in}{0.526991in}}%
\pgfpathlineto{\pgfqpoint{1.574343in}{0.543617in}}%
\pgfpathlineto{\pgfqpoint{1.582586in}{0.563930in}}%
\pgfpathlineto{\pgfqpoint{1.590828in}{0.587826in}}%
\pgfpathlineto{\pgfqpoint{1.603193in}{0.630172in}}%
\pgfpathlineto{\pgfqpoint{1.615556in}{0.680033in}}%
\pgfpathlineto{\pgfqpoint{1.627921in}{0.737093in}}%
\pgfpathlineto{\pgfqpoint{1.640283in}{0.801051in}}%
\pgfpathlineto{\pgfqpoint{1.656768in}{0.896557in}}%
\pgfpathlineto{\pgfqpoint{1.674616in}{1.012125in}}%
\pgfpathlineto{\pgfqpoint{1.684953in}{1.074230in}}%
\pgfpathlineto{\pgfqpoint{1.693195in}{1.116028in}}%
\pgfpathlineto{\pgfqpoint{1.701438in}{1.151170in}}%
\pgfpathlineto{\pgfqpoint{1.709680in}{1.179835in}}%
\pgfpathlineto{\pgfqpoint{1.717923in}{1.202198in}}%
\pgfpathlineto{\pgfqpoint{1.726165in}{1.218432in}}%
\pgfpathlineto{\pgfqpoint{1.730285in}{1.224303in}}%
\pgfpathlineto{\pgfqpoint{1.734407in}{1.228703in}}%
\pgfpathlineto{\pgfqpoint{1.738527in}{1.231654in}}%
\pgfpathlineto{\pgfqpoint{1.742650in}{1.233175in}}%
\pgfpathlineto{\pgfqpoint{1.746770in}{1.233286in}}%
\pgfpathlineto{\pgfqpoint{1.750892in}{1.232007in}}%
\pgfpathlineto{\pgfqpoint{1.755012in}{1.229357in}}%
\pgfpathlineto{\pgfqpoint{1.759135in}{1.225354in}}%
\pgfpathlineto{\pgfqpoint{1.763255in}{1.220019in}}%
\pgfpathlineto{\pgfqpoint{1.771497in}{1.205422in}}%
\pgfpathlineto{\pgfqpoint{1.779739in}{1.185711in}}%
\pgfpathlineto{\pgfqpoint{1.787982in}{1.161031in}}%
\pgfpathlineto{\pgfqpoint{1.796224in}{1.131518in}}%
\pgfpathlineto{\pgfqpoint{1.804467in}{1.097309in}}%
\pgfpathlineto{\pgfqpoint{1.816832in}{1.037478in}}%
\pgfpathlineto{\pgfqpoint{1.829194in}{0.967810in}}%
\pgfpathlineto{\pgfqpoint{1.841559in}{0.888719in}}%
\pgfpathlineto{\pgfqpoint{1.862828in}{0.744894in}}%
\pgfpathlineto{\pgfqpoint{1.875193in}{0.676058in}}%
\pgfpathlineto{\pgfqpoint{1.883436in}{0.636876in}}%
\pgfpathlineto{\pgfqpoint{1.891678in}{0.602894in}}%
\pgfpathlineto{\pgfqpoint{1.899921in}{0.573974in}}%
\pgfpathlineto{\pgfqpoint{1.908163in}{0.549979in}}%
\pgfpathlineto{\pgfqpoint{1.916405in}{0.530777in}}%
\pgfpathlineto{\pgfqpoint{1.924648in}{0.516238in}}%
\pgfpathlineto{\pgfqpoint{1.932890in}{0.506237in}}%
\pgfpathlineto{\pgfqpoint{1.937010in}{0.502899in}}%
\pgfpathlineto{\pgfqpoint{1.941133in}{0.500650in}}%
\pgfpathlineto{\pgfqpoint{1.945253in}{0.499475in}}%
\pgfpathlineto{\pgfqpoint{1.949375in}{0.499359in}}%
\pgfpathlineto{\pgfqpoint{1.953495in}{0.500287in}}%
\pgfpathlineto{\pgfqpoint{1.957618in}{0.502244in}}%
\pgfpathlineto{\pgfqpoint{1.961737in}{0.505218in}}%
\pgfpathlineto{\pgfqpoint{1.969980in}{0.514157in}}%
\pgfpathlineto{\pgfqpoint{1.978222in}{0.526994in}}%
\pgfpathlineto{\pgfqpoint{1.986465in}{0.543621in}}%
\pgfpathlineto{\pgfqpoint{1.994707in}{0.563933in}}%
\pgfpathlineto{\pgfqpoint{2.002950in}{0.587830in}}%
\pgfpathlineto{\pgfqpoint{2.015315in}{0.630176in}}%
\pgfpathlineto{\pgfqpoint{2.027677in}{0.680036in}}%
\pgfpathlineto{\pgfqpoint{2.040042in}{0.737096in}}%
\pgfpathlineto{\pgfqpoint{2.052404in}{0.801055in}}%
\pgfpathlineto{\pgfqpoint{2.068889in}{0.896560in}}%
\pgfpathlineto{\pgfqpoint{2.086737in}{1.012129in}}%
\pgfpathlineto{\pgfqpoint{2.097074in}{1.074233in}}%
\pgfpathlineto{\pgfqpoint{2.105317in}{1.116032in}}%
\pgfpathlineto{\pgfqpoint{2.113559in}{1.151174in}}%
\pgfpathlineto{\pgfqpoint{2.121801in}{1.179838in}}%
\pgfpathlineto{\pgfqpoint{2.130044in}{1.202202in}}%
\pgfpathlineto{\pgfqpoint{2.138286in}{1.218435in}}%
\pgfpathlineto{\pgfqpoint{2.142406in}{1.224306in}}%
\pgfpathlineto{\pgfqpoint{2.146529in}{1.228707in}}%
\pgfpathlineto{\pgfqpoint{2.150649in}{1.231657in}}%
\pgfpathlineto{\pgfqpoint{2.154771in}{1.233179in}}%
\pgfpathlineto{\pgfqpoint{2.158891in}{1.233290in}}%
\pgfpathlineto{\pgfqpoint{2.163013in}{1.232011in}}%
\pgfpathlineto{\pgfqpoint{2.167133in}{1.229360in}}%
\pgfpathlineto{\pgfqpoint{2.171256in}{1.225358in}}%
\pgfpathlineto{\pgfqpoint{2.175376in}{1.220022in}}%
\pgfpathlineto{\pgfqpoint{2.183618in}{1.205425in}}%
\pgfpathlineto{\pgfqpoint{2.191861in}{1.185715in}}%
\pgfpathlineto{\pgfqpoint{2.200103in}{1.161034in}}%
\pgfpathlineto{\pgfqpoint{2.208345in}{1.131521in}}%
\pgfpathlineto{\pgfqpoint{2.216588in}{1.097313in}}%
\pgfpathlineto{\pgfqpoint{2.228953in}{1.037482in}}%
\pgfpathlineto{\pgfqpoint{2.241315in}{0.967813in}}%
\pgfpathlineto{\pgfqpoint{2.253680in}{0.888723in}}%
\pgfpathlineto{\pgfqpoint{2.274950in}{0.744897in}}%
\pgfpathlineto{\pgfqpoint{2.287315in}{0.676061in}}%
\pgfpathlineto{\pgfqpoint{2.295557in}{0.636879in}}%
\pgfpathlineto{\pgfqpoint{2.303799in}{0.602898in}}%
\pgfpathlineto{\pgfqpoint{2.312042in}{0.573978in}}%
\pgfpathlineto{\pgfqpoint{2.320284in}{0.549982in}}%
\pgfpathlineto{\pgfqpoint{2.328527in}{0.530780in}}%
\pgfpathlineto{\pgfqpoint{2.336769in}{0.516242in}}%
\pgfpathlineto{\pgfqpoint{2.345011in}{0.506240in}}%
\pgfpathlineto{\pgfqpoint{2.349131in}{0.502903in}}%
\pgfpathlineto{\pgfqpoint{2.353254in}{0.500654in}}%
\pgfpathlineto{\pgfqpoint{2.357374in}{0.499479in}}%
\pgfpathlineto{\pgfqpoint{2.361496in}{0.499362in}}%
\pgfpathlineto{\pgfqpoint{2.365616in}{0.500290in}}%
\pgfpathlineto{\pgfqpoint{2.369739in}{0.502248in}}%
\pgfpathlineto{\pgfqpoint{2.373859in}{0.505222in}}%
\pgfpathlineto{\pgfqpoint{2.382101in}{0.514161in}}%
\pgfpathlineto{\pgfqpoint{2.390343in}{0.526997in}}%
\pgfpathlineto{\pgfqpoint{2.398586in}{0.543624in}}%
\pgfpathlineto{\pgfqpoint{2.406828in}{0.563937in}}%
\pgfpathlineto{\pgfqpoint{2.415071in}{0.587833in}}%
\pgfpathlineto{\pgfqpoint{2.427436in}{0.630179in}}%
\pgfpathlineto{\pgfqpoint{2.439798in}{0.680039in}}%
\pgfpathlineto{\pgfqpoint{2.452163in}{0.737100in}}%
\pgfpathlineto{\pgfqpoint{2.464525in}{0.801058in}}%
\pgfpathlineto{\pgfqpoint{2.481010in}{0.896564in}}%
\pgfpathlineto{\pgfqpoint{2.498859in}{1.012132in}}%
\pgfpathlineto{\pgfqpoint{2.509195in}{1.074237in}}%
\pgfpathlineto{\pgfqpoint{2.517438in}{1.116035in}}%
\pgfpathlineto{\pgfqpoint{2.525680in}{1.151177in}}%
\pgfpathlineto{\pgfqpoint{2.533923in}{1.179842in}}%
\pgfpathlineto{\pgfqpoint{2.542165in}{1.202205in}}%
\pgfpathlineto{\pgfqpoint{2.550407in}{1.218439in}}%
\pgfpathlineto{\pgfqpoint{2.554527in}{1.224309in}}%
\pgfpathlineto{\pgfqpoint{2.558650in}{1.228710in}}%
\pgfpathlineto{\pgfqpoint{2.562770in}{1.231661in}}%
\pgfpathlineto{\pgfqpoint{2.566892in}{1.233182in}}%
\pgfpathlineto{\pgfqpoint{2.571012in}{1.233293in}}%
\pgfpathlineto{\pgfqpoint{2.575135in}{1.232014in}}%
\pgfpathlineto{\pgfqpoint{2.579255in}{1.229364in}}%
\pgfpathlineto{\pgfqpoint{2.583377in}{1.225361in}}%
\pgfpathlineto{\pgfqpoint{2.587497in}{1.220026in}}%
\pgfpathlineto{\pgfqpoint{2.595739in}{1.205428in}}%
\pgfpathlineto{\pgfqpoint{2.603982in}{1.185718in}}%
\pgfpathlineto{\pgfqpoint{2.612224in}{1.161037in}}%
\pgfpathlineto{\pgfqpoint{2.620467in}{1.131525in}}%
\pgfpathlineto{\pgfqpoint{2.628709in}{1.097316in}}%
\pgfpathlineto{\pgfqpoint{2.641074in}{1.037485in}}%
\pgfpathlineto{\pgfqpoint{2.653436in}{0.967816in}}%
\pgfpathlineto{\pgfqpoint{2.665801in}{0.888726in}}%
\pgfpathlineto{\pgfqpoint{2.687071in}{0.744901in}}%
\pgfpathlineto{\pgfqpoint{2.699436in}{0.676064in}}%
\pgfpathlineto{\pgfqpoint{2.707678in}{0.636883in}}%
\pgfpathlineto{\pgfqpoint{2.715921in}{0.602901in}}%
\pgfpathlineto{\pgfqpoint{2.724163in}{0.573981in}}%
\pgfpathlineto{\pgfqpoint{2.732405in}{0.549986in}}%
\pgfpathlineto{\pgfqpoint{2.740648in}{0.530784in}}%
\pgfpathlineto{\pgfqpoint{2.748890in}{0.516245in}}%
\pgfpathlineto{\pgfqpoint{2.757133in}{0.506244in}}%
\pgfpathlineto{\pgfqpoint{2.761253in}{0.502906in}}%
\pgfpathlineto{\pgfqpoint{2.765375in}{0.500657in}}%
\pgfpathlineto{\pgfqpoint{2.769495in}{0.499482in}}%
\pgfpathlineto{\pgfqpoint{2.773618in}{0.499366in}}%
\pgfpathlineto{\pgfqpoint{2.777737in}{0.500293in}}%
\pgfpathlineto{\pgfqpoint{2.781860in}{0.502251in}}%
\pgfpathlineto{\pgfqpoint{2.785980in}{0.505225in}}%
\pgfpathlineto{\pgfqpoint{2.794222in}{0.514164in}}%
\pgfpathlineto{\pgfqpoint{2.802465in}{0.527001in}}%
\pgfpathlineto{\pgfqpoint{2.810707in}{0.543628in}}%
\pgfpathlineto{\pgfqpoint{2.818950in}{0.563940in}}%
\pgfpathlineto{\pgfqpoint{2.827192in}{0.587836in}}%
\pgfpathlineto{\pgfqpoint{2.839557in}{0.630183in}}%
\pgfpathlineto{\pgfqpoint{2.851919in}{0.680043in}}%
\pgfpathlineto{\pgfqpoint{2.864284in}{0.737103in}}%
\pgfpathlineto{\pgfqpoint{2.876647in}{0.801061in}}%
\pgfpathlineto{\pgfqpoint{2.893131in}{0.896567in}}%
\pgfpathlineto{\pgfqpoint{2.910980in}{1.012135in}}%
\pgfpathlineto{\pgfqpoint{2.921317in}{1.074240in}}%
\pgfpathlineto{\pgfqpoint{2.929559in}{1.116039in}}%
\pgfpathlineto{\pgfqpoint{2.937801in}{1.151180in}}%
\pgfpathlineto{\pgfqpoint{2.946044in}{1.179845in}}%
\pgfpathlineto{\pgfqpoint{2.954286in}{1.202208in}}%
\pgfpathlineto{\pgfqpoint{2.962529in}{1.218442in}}%
\pgfpathlineto{\pgfqpoint{2.966649in}{1.224313in}}%
\pgfpathlineto{\pgfqpoint{2.970771in}{1.228713in}}%
\pgfpathlineto{\pgfqpoint{2.974891in}{1.231664in}}%
\pgfpathlineto{\pgfqpoint{2.979013in}{1.233185in}}%
\pgfpathlineto{\pgfqpoint{2.983133in}{1.233297in}}%
\pgfpathlineto{\pgfqpoint{2.987256in}{1.232017in}}%
\pgfpathlineto{\pgfqpoint{2.991376in}{1.229367in}}%
\pgfpathlineto{\pgfqpoint{2.995498in}{1.225365in}}%
\pgfpathlineto{\pgfqpoint{2.999618in}{1.220029in}}%
\pgfpathlineto{\pgfqpoint{3.007861in}{1.205432in}}%
\pgfpathlineto{\pgfqpoint{3.016103in}{1.185722in}}%
\pgfpathlineto{\pgfqpoint{3.024345in}{1.161041in}}%
\pgfpathlineto{\pgfqpoint{3.032588in}{1.131528in}}%
\pgfpathlineto{\pgfqpoint{3.040830in}{1.097319in}}%
\pgfpathlineto{\pgfqpoint{3.053195in}{1.037489in}}%
\pgfpathlineto{\pgfqpoint{3.065558in}{0.967820in}}%
\pgfpathlineto{\pgfqpoint{3.077923in}{0.888729in}}%
\pgfpathlineto{\pgfqpoint{3.099192in}{0.744904in}}%
\pgfpathlineto{\pgfqpoint{3.111557in}{0.676068in}}%
\pgfpathlineto{\pgfqpoint{3.119799in}{0.636886in}}%
\pgfpathlineto{\pgfqpoint{3.128042in}{0.602905in}}%
\pgfpathlineto{\pgfqpoint{3.136284in}{0.573984in}}%
\pgfpathlineto{\pgfqpoint{3.144527in}{0.549989in}}%
\pgfpathlineto{\pgfqpoint{3.152769in}{0.530787in}}%
\pgfpathlineto{\pgfqpoint{3.161011in}{0.516248in}}%
\pgfpathlineto{\pgfqpoint{3.169254in}{0.506247in}}%
\pgfpathlineto{\pgfqpoint{3.173374in}{0.502910in}}%
\pgfpathlineto{\pgfqpoint{3.177496in}{0.500661in}}%
\pgfpathlineto{\pgfqpoint{3.181616in}{0.499486in}}%
\pgfpathlineto{\pgfqpoint{3.185739in}{0.499369in}}%
\pgfpathlineto{\pgfqpoint{3.189859in}{0.500297in}}%
\pgfpathlineto{\pgfqpoint{3.193981in}{0.502255in}}%
\pgfpathlineto{\pgfqpoint{3.198101in}{0.505229in}}%
\pgfpathlineto{\pgfqpoint{3.206343in}{0.514168in}}%
\pgfpathlineto{\pgfqpoint{3.214586in}{0.527004in}}%
\pgfpathlineto{\pgfqpoint{3.222828in}{0.543631in}}%
\pgfpathlineto{\pgfqpoint{3.231071in}{0.563944in}}%
\pgfpathlineto{\pgfqpoint{3.239313in}{0.587840in}}%
\pgfpathlineto{\pgfqpoint{3.251678in}{0.630186in}}%
\pgfpathlineto{\pgfqpoint{3.264040in}{0.680046in}}%
\pgfpathlineto{\pgfqpoint{3.276405in}{0.737107in}}%
\pgfpathlineto{\pgfqpoint{3.288768in}{0.801065in}}%
\pgfpathlineto{\pgfqpoint{3.305253in}{0.896571in}}%
\pgfpathlineto{\pgfqpoint{3.320983in}{0.998048in}}%
\pgfpathlineto{\pgfqpoint{3.320983in}{0.998048in}}%
\pgfusepath{stroke}%
\end{pgfscope}%
\begin{pgfscope}%
\pgfsetrectcap%
\pgfsetmiterjoin%
\pgfsetlinewidth{0.803000pt}%
\definecolor{currentstroke}{rgb}{0.000000,0.000000,0.000000}%
\pgfsetstrokecolor{currentstroke}%
\pgfsetdash{}{0pt}%
\pgfpathmoveto{\pgfqpoint{0.724619in}{0.462654in}}%
\pgfpathlineto{\pgfqpoint{0.724619in}{1.269994in}}%
\pgfusepath{stroke}%
\end{pgfscope}%
\begin{pgfscope}%
\pgfsetrectcap%
\pgfsetmiterjoin%
\pgfsetlinewidth{0.803000pt}%
\definecolor{currentstroke}{rgb}{0.000000,0.000000,0.000000}%
\pgfsetstrokecolor{currentstroke}%
\pgfsetdash{}{0pt}%
\pgfpathmoveto{\pgfqpoint{3.444619in}{0.462654in}}%
\pgfpathlineto{\pgfqpoint{3.444619in}{1.269994in}}%
\pgfusepath{stroke}%
\end{pgfscope}%
\begin{pgfscope}%
\pgfsetrectcap%
\pgfsetmiterjoin%
\pgfsetlinewidth{0.803000pt}%
\definecolor{currentstroke}{rgb}{0.000000,0.000000,0.000000}%
\pgfsetstrokecolor{currentstroke}%
\pgfsetdash{}{0pt}%
\pgfpathmoveto{\pgfqpoint{0.724619in}{0.462654in}}%
\pgfpathlineto{\pgfqpoint{3.444619in}{0.462654in}}%
\pgfusepath{stroke}%
\end{pgfscope}%
\begin{pgfscope}%
\pgfsetrectcap%
\pgfsetmiterjoin%
\pgfsetlinewidth{0.803000pt}%
\definecolor{currentstroke}{rgb}{0.000000,0.000000,0.000000}%
\pgfsetstrokecolor{currentstroke}%
\pgfsetdash{}{0pt}%
\pgfpathmoveto{\pgfqpoint{0.724619in}{1.269994in}}%
\pgfpathlineto{\pgfqpoint{3.444619in}{1.269994in}}%
\pgfusepath{stroke}%
\end{pgfscope}%
\begin{pgfscope}%
\definecolor{textcolor}{rgb}{0.000000,0.000000,0.000000}%
\pgfsetstrokecolor{textcolor}%
\pgfsetfillcolor{textcolor}%
\pgftext[x=3.254219in,y=1.213480in,left,top]{\color{textcolor}{\rmfamily\fontsize{8.000000}{9.600000}\selectfont\catcode`\^=\active\def^{\ifmmode\sp\else\^{}\fi}\catcode`\%=\active\def%{\%}(b)}}%
\end{pgfscope}%
\end{pgfpicture}%
\makeatother%
\endgroup%

	\caption{(a) corrente de entrada e (b) corrente de saída do circuito Cuk.}
	\label{fig:ex3_3}
\end{figure}